\documentclass[12pt]{article}
\usepackage[utf8]{inputenc}
\usepackage{amsmath, amssymb, amsthm}
\usepackage{mathtools}
\usepackage{geometry}
\usepackage{hyperref}
\usepackage{bm}
\usepackage{booktabs}
\usepackage{setspace}
\usepackage{xcolor}
\usepackage{graphicx}
\usepackage{enumitem}
\usepackage{caption}
\usepackage{subcaption}
\usepackage{longtable}
\usepackage{float}
\usepackage[english]{babel}
\onehalfspacing
\geometry{margin=1.2in}
\hypersetup{colorlinks=true, linkcolor=blue!50!black, citecolor=blue!50!black, urlcolor=blue!50!black}

\newtheorem{lemma}{Lemma}
\newtheorem{proposition}{Proposition}
\newtheorem{corollary}{Corollary}
\newtheorem{definition}{Definition}
\newtheorem{assumption}{Assumption}
\newtheorem{remark}{Remark}

\title{\Large Networks, Phillips Curves, and Exchange Rate Policy: \\[4pt]
A Small Open Economy Extension of Rubbo (2024)}
\author{Sugarkhuu \thanks{I thank Christian, Benny, and Adam for extensive comments on an earlier
presentation of this material, which shaped much of the robustness analysis in Sections
\ref{sec:results} and \ref{sec:robustness}. All errors are my own.}}
\date{\today}

\begin{document}
\maketitle

\begin{abstract}
\noindent I extend Rubbo's (2024) closed-economy theory of Phillips curves in production networks
to a small open economy (SOE) with an exchange-rate channel, and ask whether her ``divine
coincidence'' result --- the existence of a single inflation index whose targeting simultaneously
closes the output gap --- survives when the exchange rate becomes a second, independent cost-push
channel alongside sectoral productivity. I show theoretically that it generically does not: the
divine-coincidence index is the unique left null vector of a cost-push matrix that must additionally
annihilate the exchange-rate pass-through vector, two conditions that a single vector can jointly
satisfy only in the knife-edge case where pass-through is proportional to labor intensity. I take the
model to three real national input--output calibrations (Chile, South Korea, Czechia) and compare
float, peg, and managed-float exchange-rate regimes. Managed float dominates in every calibration;
peg is the worst regime, and its loss is driven overwhelmingly (58--83\% depending on the calibration
of the debt-elastic risk-premium coefficient) by a purely financial risk-premium/uncovered-interest-parity
shock rather than by the terms-of-trade channel emphasized in the classical Mundell--Fleming and
Gal\'i--Monacelli literature. An extensive robustness campaign --- roughly 350 additional model
solves spanning network density, network topology, uniform and sector-specific price rigidity, the
interaction between the risk-premium shock and its closing-device elasticity, and alternative
policy-rule designs --- confirms the ranking survives essentially everywhere, while also uncovering
that (i) a sufficiently aggressive plain inflation-targeting rule can rival the managed float's
performance without any explicit exchange-rate term, and (ii) once policy is aggressive, targeting
domestic (PPI) inflation outperforms targeting the theoretically ``correct'' divine-coincidence
index. I discuss the paper's relation to the contemporaneous and independent Qiu, Wang, Xu and
Zanetti (2026) result, and lay out the concrete steps --- a fully disaggregated multi-sector
calibration and an estimated (rather than assumed) risk-premium process chief among them --- required
before the quantitative ranking can be treated as more than suggestive.
\end{abstract}

\tableofcontents
\newpage

%======================================================================
\section{Introduction}
%======================================================================

Modern economies are not single-good economies. Firms buy inputs from other firms, and the resulting
input--output (IO) network shapes how shocks propagate and how monetary policy should be conducted.
Rubbo (2024) formalizes this insight for a closed economy: in an $N$-sector New Keynesian model with
an arbitrary IO matrix and Calvo price-setting frequencies that differ across sectors, there exists a
unique inflation index --- weighted by each sector's Domar (total sales) weight and by its price
stickiness --- whose targeting delivers the ``divine coincidence'': zero cost-push inflation and a
simultaneously closed output gap. This is a striking generalization of the one-sector New Keynesian
benchmark, in which any targeting rule that stabilizes marginal cost also stabilizes the output gap.
Rubbo's result says that this coincidence survives the introduction of arbitrarily rich production
networks and heterogeneous price stickiness, provided the central bank targets the \emph{right}
index rather than a naive aggregate like CPI inflation.

This paper asks whether the divine coincidence survives a second, empirically central feature of
essentially every economy outside the largest few: openness. A small open economy (SOE) imports
intermediate inputs, its firms' marginal costs depend on the nominal exchange rate, and its central
bank must additionally decide how to respond to exchange-rate movements --- float, peg, or some
managed combination of the two. I show that adding this single, minimal ingredient --- an import
cost share in every sector's marginal cost, priced at the law-of-one-price exchange rate --- breaks
the divine coincidence generically. The mechanism is transparent: the exchange rate is a \emph{second}
cost-push channel, exactly analogous to the sectoral productivity shocks that generate the first one.
Rubbo's unique index is pinned down by requiring it to annihilate the productivity cost-push matrix;
requiring it to \emph{also} annihilate the exchange-rate pass-through vector is a second, generically
independent linear restriction on the same weighting vector, which a vector of finite dimension can
satisfy only in a non-generic, measure-zero case. Divine coincidence in the open economy is a
knife-edge property, not a generic one.

This theoretical result immediately raises a quantitative question with obvious policy relevance:
if no simple inflation index can costlessly close the output gap under a float, how much welfare is
actually at stake across the natural menu of exchange-rate regimes --- float, hard peg, and managed
float --- and through which shocks? I calibrate the model to real, disaggregated national
input--output data for three small open economies (Chile, South Korea, and Czechia) and solve it in
Dynare under each regime. The headline finding is that a managed float --- a conventional Taylor rule
augmented with a modest, positive response to the exchange rate --- dominates both corners in every
calibration, and that a hard peg is by a wide margin the worst regime. This ranking replicates the
classical Mundell--Fleming/Gal\'i-Monacelli conclusion that flexible exchange rates are stabilizing,
but the \emph{mechanism} is different and, I would argue, more interesting: decomposing each regime's
welfare loss by shock shows that the peg's dominance is driven overwhelmingly by a purely financial
shock --- a debt-elastic risk premium on the economy's net foreign asset position, in the tradition
of Schmitt-Grohé and Uribe's (2003) closing device --- and not by the terms-of-trade channel that is
the star of the classical open-economy literature. This is an uncomfortable result in one specific
sense that I confront directly rather than hide: since the risk-premium shock's standard deviation is
a modeling choice (set equal to every other shock's, for want of an external target), the entire
``peg is dominated'' headline finding is only as solid as that one calibration choice. I therefore
subject it, and the closely related debt-elasticity coefficient governing how strongly the risk
premium reacts to the net foreign asset position, to an extensive battery of robustness checks ---
described below --- rather than treating the headline number as settled.

The theoretical contribution of this paper sits between two literatures that, until very recently,
had not been brought together. The first is the closed-economy production-network literature
(Rubbo, 2024; La'O and Tahbaz-Salehi, 2022; Pasten, Schoenle and Weber, 2020; Baqaee and Farhi, 2019,
2020), which shows that IO linkages and heterogeneous price stickiness matter enormously for how
monetary policy should be conducted, but is silent on the exchange rate because it studies closed
economies. The second is the open-economy New Keynesian literature descending from Gal\'i and
Monacelli (2005), Obstfeld and Rogoff (2000), and Corsetti, Dedola and Leduc (2010), which analyzes
exchange-rate regimes and optimal monetary policy in richly specified open economies, but --- with
rare exceptions --- collapses the entire domestic supply side to a single sector, so that it cannot
speak to which \emph{sectors} bear the cost of a given regime or how network structure shapes the
regime ranking. A single contemporaneous paper, Qiu, Wang, Xu and Zanetti (2026), independently
derives an open-economy generalization of Rubbo's index and reaches the same qualitative
divine-coincidence-breakdown conclusion through a structurally different route (a static
Gal\'i--Monacelli trade block with expenditure-switching and profit channels, rather than this
paper's dynamic import-cost-share pass-through vector). Their model, notably, has no financial
account, no uncovered interest parity, and no exchange-rate-\emph{regime} choice --- their own stated
top-priority extension is exactly the machinery (debt-elastic risk premium, near-unit-root net
foreign assets) that this paper builds around, and that turns out to drive the entire quantitative
result. Section \ref{sec:litreview} discusses this relationship, and the surrounding literature, in
detail.

Given how much of the headline result rests on two specific, judgment-call calibration parameters
(the risk-premium shock's standard deviation and its debt-elasticity coefficient), a substantial part
of this paper --- comprising roughly 350 additional Dynare solves beyond the three headline
calibrations --- is devoted to systematically stress-testing it. I vary: the density and the
\emph{shape} of the domestic production network (a stylized triangular network scaled continuously
from zero to double density, plus two alternative topologies --- a hub-and-spoke network with
Manufacturing as the central node, and a one-directional supply chain); uniform and Services-specific
price rigidity; the debt-elasticity coefficient of the risk-premium process, on its own, jointly with
network density, and jointly with the risk-premium shock's own volatility; and the exact design of the
monetary policy rule, including strict domestic- and consumer-price-inflation targeting and a
coarse grid search over the Taylor rule's own coefficients. The central qualitative conclusion ---
Managed float dominates Float, which dominates Peg, in every single one of these variations --- is
remarkably stable. Two secondary findings from this exercise are, I think, worth flagging on their own
terms: first, that a sufficiently aggressive plain inflation-targeting rule, with no explicit
exchange-rate response at all, can rival the managed float's performance; and second, that once
policy is that aggressive, targeting domestic (producer) price inflation outperforms targeting the
theoretically privileged divine-coincidence index, reversing the ranking found under more moderate
policy coefficients. Both nuance, without overturning, the paper's central welfare ranking, and both
are flagged explicitly as preliminary rather than settled, since neither has yet been checked against
a fully re-optimized managed-float rule at the same level of policy aggressiveness.

The remainder of the paper proceeds as follows. Section \ref{sec:litreview} reviews the related
literature in three strands: closed-economy production networks, open-economy monetary policy, and
the empirical exchange-rate pass-through literature that disciplines the model's key structural
parameters. Section \ref{sec:model} lays out the small open economy model. Section
\ref{sec:theory} derives the modified sectoral Phillips curve, states and proves the divine-coincidence
breakdown result, and derives the welfare function used throughout the quantitative analysis.
Section \ref{sec:data} describes the data and calibration strategy for the three country
calibrations. Section \ref{sec:results} presents the headline quantitative results: the welfare
ranking, its shock decomposition, a mechanism analysis of why the stickiest, least import-exposed
sector (``Services'') nonetheless bears the largest share of the welfare cost, and cross-country
robustness. Section \ref{sec:robustness} presents the full robustness campaign. Section
\ref{sec:discussion} discusses policy implications and the paper's remaining limitations. Section
\ref{sec:conclusion} concludes. Appendix \ref{app:proofs} collects the detailed, step-by-step proofs
of all lemmas and propositions stated in the main text.

%======================================================================
\section{Literature Review}
\label{sec:litreview}
%======================================================================

This paper sits at the intersection of three literatures: production networks and monetary policy in
closed economies, exchange-rate policy and the divine coincidence in (single-sector) open economies,
and the empirical literature on how production networks shape international pass-through. I discuss
each in turn, before turning to the single closest paper in the literature.

\subsection{Production Networks and Monetary Policy in Closed Economies}

Rubbo (2024) is the direct theoretical parent of this paper. In an $N$-sector economy with Calvo
price-setting, heterogeneous sectoral stickiness $\hat\delta_i$, and an arbitrary IO matrix $\Omega$,
she shows three results this paper builds on directly. First, natural (flexible-price) output equals
Domar-weighted total factor productivity (a version of Hulten's, 1978, theorem applied to the social
planner's problem). Second, sectoral inflation rates satisfy a vector Phillips curve
$\bm\pi_t = \rho(I-\mathcal V)\mathbb E\bm\pi_{t+1} + \mathcal B(\gamma+\varphi)\tilde y_t -
\mathcal V\bm\chi_t$, in which the cost-push matrix $\mathcal V$ has rows that sum to zero --- only
\emph{relative}, not uniform, productivity shocks generate cost-push inflation --- and is built from
the rigidity-adjusted Leontief inverse $(I-\Omega\Delta)^{-1}$, so that stickier upstream sectors
dampen the propagation of shocks by absorbing them into markups rather than passing them through as
price changes. Third, and most importantly for this paper, there is a unique inflation index
$\pi_t^{DC}=\sum_i w_i^{DC}\pi_{it}$, with weights proportional to $\lambda_i(1-\hat\delta_i)/\hat\delta_i$
(Domar weight times a stickiness multiplier), whose own law of motion has no cost-push term at all:
targeting it in a standard Taylor-type rule simultaneously delivers a zero output gap. Rubbo further
shows that a naive CPI-targeting rule, which ignores this network structure, can generate welfare
losses of 2.9--3.8\% of GDP relative to the divine-coincidence optimum even in the closed economy ---
motivation for taking network structure seriously as a first-order input into monetary policy design,
not a second-order refinement.

La'O and Tahbaz-Salehi (2022) reach a closely related conclusion from a different direction: using a
non-parametric, sufficient-statistics approach (rather than Rubbo's log-linear approximation), they
show that optimal monetary policy in a production-network economy depends only on sectoral sales
shares, demand elasticities, and the degree of strategic complementarity in pricing, and that the
welfare-relevant inflation target weights sectors by Domar weight and inverse markup sensitivity ---
not by consumption share, as a naive CPI target would. Pasten, Schoenle and Weber (2020) provide
the quantitative counterpart, calibrating a multi-sector Calvo model to U.S. input--output and Calvo
frequency data (from Nakamura and Steinsson, 2008) and showing that network linkages amplify the real
effects of a monetary policy shock by a factor of roughly two to three relative to an equivalent
one-sector economy, with sticky, high-Domar-weight (frequently downstream, service-sector-like)
industries accounting for a disproportionate share of the amplification --- a pattern this paper's
own ``Services'' mechanism, discussed in Section \ref{sec:results}, closely parallels in the open
economy.

Underlying all of the above is the foundational network-macro literature: Hulten (1978) and Acemoglu,
Carvalho, Ozdaglar and Tahbaz-Salehi (2012) establish that idiosyncratic sectoral shocks need not
average out at the standard $1/\sqrt N$ rate once a fat-tailed Domar-weight distribution is present,
so that individually small shocks to central (``hub'') sectors can generate aggregate fluctuations.
Baqaee and Farhi (2019, 2020) generalize Hulten's first-order result to second order, showing that
the exact aggregate effect of sectoral distortions requires knowing not just Domar weights but the
full pattern of substitution elasticities across sectors; the cross-sector welfare terms
$\Phi_C,\Phi_s$ that appear in Rubbo's Proposition 3, and in this paper's welfare function
(Section \ref{sec:theory}), are precisely these Baqaee--Farhi operators. Bigio and La'O (2020) apply
the same network-aggregation logic to arbitrary wedges (taxes, markups, financial frictions) rather
than pure productivity shocks --- the natural lens through which to view this paper's exchange-rate
cost-push term, which is, formally, a heterogeneous wedge on each sector's imported-input cost.

\subsection{Exchange-Rate Policy and the Divine Coincidence in Open Economies}

The benchmark open-economy New Keynesian model is Gal\'i and Monacelli (2005): a small open economy
facing a large foreign economy, with households consuming a CES aggregate of home and foreign goods,
firms using only labor (no intermediate inputs, let alone a production network), complete
international financial markets, and the law of one price. In this environment, targeting domestic
(PPI) inflation delivers the one-sector divine coincidence exactly as in a closed economy, with the
flexible exchange rate acting as an automatic stabilizer of the terms of trade; targeting CPI
inflation instead is strictly suboptimal because it forces the central bank to partially resist
exchange-rate movements that would otherwise be efficient; and a hard peg is welfare-inferior because
it surrenders monetary independence (the Mundell trilemma) in the face of domestic productivity
shocks. Corsetti, Dedola and Leduc (2010) survey the conditions under which this benchmark divine
coincidence survives richer environments --- incomplete markets, pricing-to-market, local-currency
pricing --- and under which it breaks down; the production-network channel this paper introduces is
absent from every model they survey.

A second relevant strand concerns the financial side of the small open economy. Schmitt-Grohé and
Uribe's (2003) debt-elastic interest-rate premium is the standard closing device for models with
incomplete asset markets and a stationary net-foreign-asset position; this paper adopts it directly,
and much of the robustness analysis in Section \ref{sec:robustness} is devoted to stress-testing the
model's sensitivity to its calibration. Fanelli and Straub (2021) provide the theoretical
underpinning for treating foreign-exchange intervention as a genuine second policy instrument
alongside the interest rate, showing formally that a central bank facing a UIP risk premium should
lean against the component of exchange-rate movements that generates inefficient real effects while
allowing the fundamentals-driven component to pass through --- exactly the logic behind this paper's
``managed float'' regime, which augments a standard Taylor rule with a response to the exchange rate.
Gopinath and Itskhoki's (2022) dominant-currency-paradigm literature is a natural extension this paper
does not pursue: the baseline model assumes producer-currency pricing throughout, and allowing a
fraction of imports to be invoiced in a third (dominant) currency would alter the exchange-rate
pass-through vector in ways left for future work. Schmitt-Grohé and Uribe (2016) show that adding
downward nominal wage rigidity to a pegged small open economy generates large, persistent involuntary
unemployment following external shocks --- since this paper's model has fully flexible wages, this
suggests the peg's welfare cost documented here is, if anything, a conservative lower bound.

\subsection{Production Networks and International Pass-Through}

A largely empirical literature establishes that production-network linkages are quantitatively
important for how prices and shocks transmit across borders, disciplining this paper's calibration of
the exchange-rate pass-through vector $\Gamma$. Auer, Levchenko and Sauré (2019) show, using
international input--output data for forty countries, that a one percent increase in a trading
partner's producer price inflation, weighted by bilateral IO exposure, raises domestic producer
price inflation by 0.15--0.25 percent above and beyond standard trade-weighted exchange-rate effects,
with upstream (supplier-side) linkages mattering more than downstream ones --- direct motivation for
this paper's $\Gamma\Delta e_t$ cost-push term, and for using disaggregated IO data rather than
aggregate trade shares to calibrate it. Amiti, Itskhoki and Konings (2014) show at the firm level
that import intensity is the key determinant of incomplete exchange-rate pass-through to export
prices (a natural hedge: depreciation raises both a firm's foreign revenue and its imported-input
costs), consistent with this paper's assumption that $\Gamma_i$ is proportional to sector $i$'s
import cost share, amplified through the network. Bems and Johnson (2017) develop
``value-added exchange rates,'' showing that standard trade-weighted real effective exchange rates
are systematically biased for economies deeply embedded in global value chains --- the empirical
counterpart of this paper's insight that a sector's \emph{effective} exposure to the exchange rate
depends on the entire chain of its suppliers' import intensities, not just its own. Nakamura and
Steinsson (2008) and Carvalho (2006) document the extensive cross-sector heterogeneity in
price-adjustment frequency that gives Rubbo's, and this paper's, $\hat\delta_i$-weighted objects their
empirical bite; without this heterogeneity, the divine-coincidence index would collapse to a simple
Domar-weighted (or, under Cobb--Douglas consumption, consumption-weighted) average and the
distinction this paper studies would be moot.

\subsection{The Closest Paper: Qiu, Wang, Xu and Zanetti (2026)}

The single closest paper in the literature is Qiu, Wang, Xu and Zanetti (2026, \emph{Journal of
Monetary Economics}), an independent and contemporaneous extension of Rubbo's framework to the open
economy. Their model combines Gal\'i--Monacelli's one-sector open-economy trade block with Rubbo's
multi-sector production network, calibrated to the World Input-Output Database (43 countries by 56
sectors). They derive an open-economy generalization of the divine-coincidence index in which the
output-gap weight on each sector decomposes into a CPI channel, an expenditure-switching channel, and
a profit channel (the latter two specific to the open economy), and show --- as this paper does,
through a structurally different mechanism --- that full divine coincidence (zero welfare loss, not
merely a zero output gap) fails even after their generalized index is targeted, because the within-
sector, across-sector, and cross-border misallocation terms are not simultaneously zeroed by
output-gap targeting alone. Calibrated to WIOD, they show output-gap targeting is close to the
constrained-optimal policy and dominates naive PPI or CPI targeting by a margin that rises with an
economy's openness.

Three features distinguish this paper from theirs, and explain why their qualitative conclusion and
this paper's are corroborating rather than redundant. First, their model is static (one period, no
dynamics or persistence), while this paper's central mechanism --- the interaction between a
persistent, near-unit-root net-foreign-asset process and a debt-elastic risk premium --- is
inherently dynamic and could not be represented in their framework. Second, their model assumes
balanced trade with no uncovered interest parity, no net foreign assets, and no financial account at
all; the exchange rate is pinned down period-by-period by a static trade-balance condition. Third,
and consequently, their model has no exchange-rate-\emph{regime} choice: money supply targets an
inflation index and the exchange rate floats implicitly throughout, so a peg-versus-float-versus-
managed-float welfare comparison, this paper's central quantitative exercise, is outside their scope.
Tellingly, the authors themselves name ``relax[ing] the assumption of financial autarky'' as their
first, top-priority item for future research --- precisely the piece of machinery this paper is built
around. This paper's contribution, then, is best read as picking up exactly where Qiu, Wang, Xu and
Zanetti's own stated research agenda leaves off, and the finding that the risk-premium/UIP channel
--- not the terms-of-trade or expenditure-switching channels their model emphasizes --- turns out to
dominate the exchange-rate-regime welfare ranking once that machinery is added.

A related, purely empirical paper is Silva (2024), who applies a production-network SOE framework to
Chilean and UK data to explain the 2021--2022 inflation surge via a CPI-elasticity decomposition
distinguishing direct from network-inherited import and export exposure. Silva's paper is positive
rather than normative --- it does not compare monetary or exchange-rate regimes --- but its Chilean
application provides a useful external check on the plausibility of this paper's own Chile-calibrated
pass-through magnitudes, and the general finding that indirect (network-inherited) exposure matters
quantitatively for inflation dynamics is the same mechanism this paper's ``Why Services?'' analysis
(Section \ref{sec:results}) documents in the welfare, rather than the positive-inflation, domain.

%======================================================================
\section{Model}
\label{sec:model}
%======================================================================

The model extends Rubbo's (2024) closed-economy $N$-sector New Keynesian production-network economy
to a small open economy with an imported-input sector, incomplete international financial markets,
and a menu of exchange-rate policy regimes. I present the environment in full; the log-linearization
and all formal results are proved in Appendix \ref{app:proofs}.

\subsection{Households}

A representative household maximizes
\begin{equation}
    \mathbb E_0\sum_{t=0}^\infty\rho^t\left[\frac{C_t^{1-\gamma}}{1-\gamma}-\frac{L_t^{1+\varphi}}{1+\varphi}\right]
\end{equation}
subject to the flow budget constraint
\begin{equation}
    P_t^C C_t + B_{t+1} + \mathcal E_t B_{t+1}^* \le W_tL_t + \Pi_t - T_t + (1+i_t)B_t
    + \mathcal E_t(1+i_t^*)\,\mathrm{RP}_t\,B_t^*
\end{equation}
where $C_t$ is a consumption aggregate of domestic and foreign varieties, $B_t$ is a domestic bond
(in zero net supply in equilibrium, but included explicitly rather than assumed away),
$B_t^*$ is the economy's net foreign asset position (in units of foreign currency), $\mathcal E_t$ is
the nominal exchange rate (domestic currency per unit of foreign currency, so that an increase is a
depreciation), and $\mathrm{RP}_t$ is a time-varying risk-premium wedge on foreign-currency returns,
discussed below. Consumption $C_t$ is a CES aggregate of a domestic consumption bundle $C_t^H$ (itself
Cobb--Douglas across the $N$ domestic sectors with shares $\bm\beta^H$) and a foreign consumption
composite $C_t^F$ (a single Armington import aggregate, not sector-indexed),
\begin{equation}
    C_t = \left[\omega^{1/\eta}(C_t^H)^{(\eta-1)/\eta} + (1-\omega)^{1/\eta}(C_t^F)^{(\eta-1)/\eta}\right]^{\eta/(\eta-1)}
\end{equation}
with $\eta$ the (CES) elasticity of substitution between home and foreign consumption bundles and
$\omega$ the home bias parameter. The household's intratemporal and intertemporal optimality
conditions are the standard labor-supply condition $W_t/P_t^C = C_t^\gamma L_t^\varphi$ and Euler
equation, exactly as in the closed economy; the presence of the foreign bond additionally delivers an
uncovered interest parity (UIP) condition, given below.

\subsection{Firms, Technology, and the Production Network}

There are $N$ domestic industries (in the quantitative model, $N=3$: Resource, Manufacturing, and
Services), each populated by a unit mass of monopolistically competitive firms. A firm $f$ in sector
$i$ produces using labor, domestic intermediate inputs from every sector $j$, and imported inputs,
\begin{equation}
    Y_{ift} = A_{it}\,L_{ift}^{\alpha_i}\prod_{j=1}^N X_{ijft}^{\Omega^H_{ij}}\,M_{ift}^{\Omega^F_i},
    \qquad \alpha_i + \sum_j\Omega^H_{ij} + \Omega^F_i = 1,
\end{equation}
Cobb--Douglas across all of labor, the $N$ domestic inputs, and the imported-input composite $M_{ift}$
(priced, under the law of one price and producer-currency pricing, at $P_t^F = \mathcal E_t P_t^*$).
$\Omega^H$ is the $N\times N$ domestic input--output matrix exactly as in Rubbo (2024); $\bm\Omega^F$
is the new object this paper introduces, the vector of sector-level imported-input cost shares.
Within-sector demand and the Calvo price-setting mechanism are unchanged from the closed economy: a
fraction $\delta_i$ of firms in sector $i$ reset their price each period, generating, after
log-linearization, an effective Calvo parameter $\hat\delta_i \equiv
\delta_i(1-\rho(1-\delta_i))/(1-\rho\delta_i(1-\delta_i))$ exactly as in Rubbo (2024). Cost
minimization gives sector $i$'s log-linearized marginal cost as
\begin{equation}
    mc_{it} = \alpha_i w_t + \sum_j\Omega^H_{ij}\,p_{jt} + \Omega^F_i\,(p_t^* + e_t) - \log A_{it}
    \label{eq:mc_main}
\end{equation}
the only modification to Rubbo's marginal cost equation being the new term $\Omega^F_i(p_t^*+e_t)$:
the nominal exchange rate $e_t\equiv\log\mathcal E_t$ enters every sector's marginal cost in
proportion to its imported-input cost share, exactly as sectoral productivity does, but with an
opposite sign and through a different cost share vector ($\bm\Omega^F$ rather than $-\log\mathbf A_t$).

\subsection{Foreign Sector, UIP, and the Risk Premium}

The economy is small: foreign variables ($P_t^*$, foreign demand $D_t^*$, the foreign interest rate
$i_t^*$) are exogenous. Uncovered interest parity, augmented with a debt-elastic risk premium in the
spirit of Schmitt-Grohé and Uribe (2003), reads
\begin{equation}
    (1+i_t) = (1+i_t^*)\,\Big(1-\Psi\big(B_t^*-\bar B^*\big)\Big)\,\mathrm{RP}_t\,\frac{\mathbb E_t\mathcal E_{t+1}}{\mathcal E_t}
    \label{eq:uip_main}
\end{equation}
where $\Psi>0$ is the debt-elasticity coefficient (governing how strongly the domestic interest rate
must compensate foreign lenders as the economy's net foreign asset position falls below its steady-
state target $\bar B^*$, the standard device restoring stationarity to an otherwise near-random-walk
net foreign asset position) and $\mathrm{RP}_t$ is an exogenous, stochastic risk-premium shock,
$\log\mathrm{RP}_t = \rho_{RP}\log\mathrm{RP}_{t-1}+\varepsilon_t^{RP}$. This risk-premium shock plays
the central quantitative role in the paper's results and is subjected to extensive robustness analysis
in Section \ref{sec:robustness}. The economy's net foreign asset position evolves according to the
balance of payments,
\begin{equation}
    B_t^* = \frac{1+i_{t-1}^*}{\Pi_t^*}\,\mathrm{RP}_{t-1}\,B_{t-1}^* + \frac{\mathrm{TB}_t}{\mathcal E_tP_t^*},
\end{equation}
where $\mathrm{TB}_t$ is the nominal trade balance (exports minus imports). Exports are driven by
foreign demand and the relative price of domestic to foreign goods, $EX_t = \kappa_{EX}D_t^*
(P_t^H/P_t^{FH})^{-\theta^*}$, with $\theta^*$ the (CES) elasticity of foreign demand for domestic
goods.

\subsection{Monetary Policy Regimes}

I compare three monetary policy regimes, the paper's central object of comparison. Under \textbf{float},
the central bank follows a conventional Taylor rule responding to the divine-coincidence index and the
output gap,
\begin{equation}
    \log(I_t/I^*) = \phi_\pi \pi_t^{DC} + \phi_y\,\tilde y_t, \tag{Float}
\end{equation}
with no direct response to the exchange rate. Under \textbf{peg}, the exchange rate is fixed,
$S_t\equiv1$, and the interest rate becomes the residual variable pinned down by the UIP equation
\eqref{eq:uip_main} --- the standard Mundell trilemma trade-off, surrendering monetary independence
in exchange for exchange-rate stability. Under \textbf{managed float}, the Taylor rule is augmented
with an explicit response to the exchange rate,
\begin{equation}
    \log(I_t/I^*) = \phi_\pi \pi_t^{DC} + \phi_y\,\tilde y_t + \phi_s\log S_t, \tag{Managed}
\end{equation}
with $\phi_s>0$ a partial FX-stabilization weight. Section \ref{sec:robustness} additionally
considers strict domestic (PPI) and consumer-price (CPI) inflation targeting, and a coarser grid
search over $(\phi_\pi,\phi_y)$ themselves, as further regime-rule variants.

\subsection{Equilibrium}

An equilibrium is a sequence of prices and quantities such that: households optimize subject to their
budget constraint; firms in every sector minimize cost and, when given the opportunity, reset prices
optimally under Calvo pricing; goods markets clear sector by sector,
$Y_{it} = C_{it}^H + \sum_j X_{jit} + EX_{it}$, where $EX_{it}$ is sector $i$'s share of aggregate
exports (allocated, in the baseline calibration, via household consumption shares $\bm\beta^H$, and
via sector-specific export intensities in the robustness extension discussed in Section
\ref{sec:robustness}); the labor market clears, $L_t=\sum_iL_{it}$; the UIP condition
\eqref{eq:uip_main} holds under float and managed float, or the peg $S_t\equiv1$ replaces it; and the
monetary authority sets the interest rate according to its regime-specific rule. GDP is a reporting
identity, $GDP_t = P_t^CC_t + P_t^H EX_t - P_t^{FH}IM_t$, and is deliberately not a market-clearing
condition.

%======================================================================
\section{Theoretical Results}
\label{sec:theory}
%======================================================================

This section states the paper's central theoretical results: the modified sectoral Phillips curve
with an exchange-rate cost-push term, the generic breakdown of the divine coincidence, and the
welfare function used throughout the quantitative analysis. Full, step-by-step proofs --- following
exactly the same log-linearization and Sherman--Morrison-identity techniques Rubbo (2024) uses in the
closed economy --- are given in Appendix \ref{app:proofs}; here I state results and give the essential
economic content of each proof.

\subsection{The Closed-Economy Benchmark}

Three results from Rubbo (2024), restated here because the open-economy results are direct
generalizations, or direct violations, of them.

\begin{lemma}[Natural output]
\label{lem:natural_output}
Natural (flexible-price) output equals Domar-weighted total factor productivity,
$y_{t,\mathrm{nat}} = \tfrac{1+\varphi}{\gamma+\varphi}\bm\lambda^\top\log\mathbf A_t$, where
$\bm\lambda^\top=\bm\beta^\top(I-\Omega^H)^{-1}$ are Domar weights.
\end{lemma}

\begin{proposition}[Sectoral Phillips curve, closed economy]
\label{prop:pc_closed}
Sectoral inflation satisfies
\begin{equation}
    \bm\pi_t = \rho(I-\mathcal V)\,\mathbb E\bm\pi_{t+1} + \mathcal B(\gamma+\varphi)\tilde y_t -
    \mathcal V\bm\chi_t,
    \label{eq:pc_closed}
\end{equation}
where $\mathcal B \equiv \Delta(I-\Omega^H\Delta)^{-1}\bm\alpha/\kappa$,
$\kappa\equiv1-\bm\beta^\top\Delta(I-\Omega^H\Delta)^{-1}\bm\alpha$, and $\mathcal V$ is a cost-push
matrix, built from the same rigidity-adjusted Leontief inverse $(I-\Omega^H\Delta)^{-1}$, whose rows
sum to zero: $\mathcal V\mathbf 1=\mathbf 0$.
\end{proposition}

\begin{proposition}[Divine coincidence, closed economy]
\label{prop:dc_closed}
There is a unique (up to normalization) inflation index $\pi_t^{DC}=\bm\phi^{DC\top}\bm\pi_t$,
$\bm\phi^{DC}_i \propto \lambda_i(1-\hat\delta_i)/\hat\delta_i$, that lies in the left null space of
$\mathcal V$ and therefore satisfies a cost-push-free Phillips curve
\begin{equation}
\pi_t^{DC}=\rho\mathbb E\pi_{t+1}^{DC} + \kappa_{DC}\tilde y_t. \label{eq:DCPC}
\end{equation}
Targeting $\pi_t^{DC}$ simultaneously
delivers a zero output gap.
\end{proposition}

Uniqueness in Proposition \ref{prop:dc_closed} follows because $\mathcal V\mathbf 1=\mathbf 0$ implies
$\mathcal V$ has rank at most $N-1$, so its left null space is generically exactly one-dimensional.
This one fact is the fulcrum on which the entire open-economy result in Section \ref{sec:dc_breakdown}
turns.

\subsection{The Modified Phillips Curve with an Exchange-Rate Channel}

Substituting the open-economy marginal cost \eqref{eq:mc_main} into the same derivation that delivers
Proposition \ref{prop:pc_closed} (Appendix \ref{app:proofs}, Steps 1--7, unchanged in every step
except the marginal cost equation itself) gives the paper's central reduced-form object.

\begin{proposition}[Sectoral Phillips curve, open economy]
\label{prop:pc_open}
Sectoral inflation satisfies
\begin{equation}
    \boxed{\bm\pi_t = \rho(I-\mathcal V)\,\mathbb E\bm\pi_{t+1} + \mathcal B(\gamma+\varphi)\tilde y_t -
    \mathcal V\bm\chi_t + \Gamma\,\Delta e_t}
    \label{eq:pc_open}
\end{equation}
where $\mathcal B$ and $\mathcal V$ are exactly as in the closed economy (Proposition
\ref{prop:pc_closed}), and the new \textbf{exchange-rate pass-through vector} is
\begin{equation}
    \Gamma \equiv \frac{\Delta(I-\Omega^H\Delta)^{-1}\bm\Omega^F}{\kappa} = \frac{\tilde A\,\bm\Omega^F}{\kappa},
    \qquad \tilde A \equiv \Delta(I-\Omega^H\Delta)^{-1}.
    \label{eq:gamma_def}
\end{equation}
\end{proposition}

Equation \eqref{eq:pc_open} is the SOE analogue of Rubbo's Proposition 2, and its economic content is
the paper's central mechanism in one line: the exchange rate enters every sector's Phillips curve
exactly like a (negative) productivity shock, weighted by the \emph{same} rigidity-adjusted network
amplifier $\tilde A=\Delta(I-\Omega^H\Delta)^{-1}$ that governs productivity pass-through, but applied
to the imported-input cost-share vector $\bm\Omega^F$ rather than to $-\log\mathbf A_t$. A sector's
exposure to the exchange rate is therefore not simply its own import share $\Omega^F_i$: it also
inherits exposure from every supplier, weighted by how quickly that supplier resets its own price.
A sector that imports nothing directly but buys heavily from an import-intensive supplier can be
substantially exposed to the exchange rate purely through the network --- the mechanism examined in
detail in Section \ref{sec:results}'s ``Why Services?'' analysis.

\subsection{Breakdown of the Divine Coincidence}
\label{sec:dc_breakdown}

\begin{proposition}[Generic breakdown of the divine coincidence]
\label{prop:dc_breakdown}
Under a float, an inflation index $\bm\phi^\top\bm\pi_t$ has no cost-push term if and only if
$\bm\phi^\top\mathcal V=\mathbf 0^\top$ \emph{and} $\bm\phi^\top\Gamma=0$. The first condition alone
pins $\bm\phi$ down uniquely (Proposition \ref{prop:dc_closed}); the second is then an
\emph{additional}, generically independent scalar restriction on the same, already-determined vector.
Except in the knife-edge case $\Gamma\propto\mathcal B$ (exchange-rate pass-through exactly
proportional to labor intensity), no inflation index exists that simultaneously eliminates both
cost-push channels. The divine coincidence is a knife-edge, not a generic, property of open production
networks.
\end{proposition}

The proof (Appendix \ref{app:proofs}) is almost immediate given Proposition \ref{prop:dc_closed}: the
closed-economy divine-coincidence weights $\bm\phi^{DC}$ are the \emph{unique} (one-dimensional) left
null vector of $\mathcal V$, so there is no freedom left to additionally satisfy
$\bm\phi^{DC\top}\Gamma=0$ unless this happens to hold already --- a single additional scalar equation
in zero remaining free parameters, generically inconsistent. The intuition is that Rubbo's closed
economy has exactly one cost-push channel (productivity), so pinning down the unique index that
annihilates it exhausts every degree of freedom available; adding the exchange rate as a
\emph{second}, generically linearly independent cost-push channel leaves no further freedom to
annihilate it as well. This is the SOE analogue of the intuition behind the Tinbergen counting
argument in Fanelli and Straub (2021): one instrument (or, here, one inflation index) cannot generally
achieve two independent targets.

\begin{remark}[Independent corroboration]
Qiu, Wang, Xu and Zanetti (2026) reach the same qualitative conclusion --- divine coincidence fails in
a multi-sector open economy --- through a structurally different route: rather than a single
import-cost-share pass-through vector $\Gamma$, they derive a three-channel open-economy output-gap
weight (CPI, expenditure-switching, profit channels) from a Gal\'i--Monacelli-style trade block with
balanced trade and export taxes, and show that output-gap targeting closes the gap but does not zero
the within-sector, across-sector, and cross-border misallocation terms simultaneously. That two
independent derivations, built on genuinely different modeling primitives, arrive at the same
qualitative conclusion is worth treating as corroboration of the underlying economics rather than as
redundant results.
\end{remark}

\subsection{Welfare}

\begin{proposition}[Welfare, open economy]
\label{prop:welfare}
To second order, the welfare loss relative to the efficient (flexible-price) allocation is
\begin{equation}
    \mathbb W = \frac12\,\mathbb E\Big[\underbrace{(\gamma+\varphi)\tilde y_t^2}_{\text{output gap}}
    + \underbrace{\sum_i\kappa_i\,\pi_{it}^2}_{\text{within-sector}}
    + \underbrace{\Phi_C(\bm\mu_t,\bm\mu_t)+\sum_s\lambda_s\Phi_s(\bm\mu_t,\bm\mu_t)}_{\text{cross-sector}}\Big],
    \qquad \kappa_i \equiv \lambda_{D,i}\,\varepsilon_i\,\frac{1-\hat\delta_i}{\hat\delta_i},
    \label{eq:welfare_main}
\end{equation}
where $\bm\mu_t$ is the vector of sectoral log markup gaps and $\Phi_C,\Phi_s$ are Baqaee--Farhi
(2019) cross-sector substitution operators evaluated at the model's consumption and production
elasticities of substitution.
\end{proposition}

The output-gap and within-sector terms are unchanged in form from the closed economy; the exchange
rate enters welfare only indirectly, through its effect on the equilibrium paths of $\tilde y_t$ and
$\bm\pi_t$ under each policy regime, plus a direct cross-sector misallocation term
$\Phi_C(\tilde{\bm\chi}_t, e_t\bm\Omega^F)$ arising because different sectors face different import
cost shocks, distorting relative prices even absent any within-sector dispersion. Because this
model's cross-sector aggregators (household consumption across domestic sectors, and each firm's
input bundle across domestic and imported inputs) are Cobb--Douglas by construction, the Allen--Uzawa
elasticity of substitution between any pair of their arguments is exactly one, and the cross-sector
terms $\Phi_C,\Phi_s$ collapse to a weighted variance of the markup-gap vector, with only the outer
consumption nest (home versus foreign consumption, $\eta=1.5$) contributing a genuine
non-unitary elasticity. Numerically (Appendix \ref{app:proofs}, Remark \ref{rem:cross_sector_numeric}),
this cross-sector term adds a modest 3.3--8.4\% on top of the output-gap-plus-within-sector total,
depending on the regime, without affecting the regime ranking; the headline welfare numbers reported
in Section \ref{sec:results} use the output-gap-plus-within-sector total (equation
\eqref{eq:welfare_main}'s first two terms) as the primary object of comparison, consistent with how
every robustness exercise in Section \ref{sec:robustness} is computed, and I flag explicitly wherever
the cross-sector term is or is not included.

%======================================================================
\section{Data and Calibration}
\label{sec:data}
%======================================================================

I calibrate the model to three small open economies with genuinely disaggregated national
input--output data: Chile, South Korea, and Czechia. Chile is the primary quantitative calibration
used throughout the paper's main results and robustness analysis; Korea and Czechia serve as
out-of-sample checks on whether the headline welfare ranking is a Chile-specific artifact or a more
general property of the model (Section \ref{sec:results}, cross-country robustness).

\subsection{Input--Output Data and Sector Aggregation}

For each country, I take the national supply-and-use (or, for Chile, Cuadro de Origen y Recursos,
CdeR) tables and collapse the native sectoral classification (12 sectors for Chile, using Banco
Central de Chile's 2018 CdeR tables at 2023 reference prices) to three broad sectors --- Resource,
Manufacturing, and Services --- chosen to span a wide range of import intensity and price stickiness
while keeping the calibrated model computationally tractable for the extensive robustness campaign in
Section \ref{sec:robustness}. From each country's IO table I compute: the domestic IO matrix
$\Omega^H$ (each sector's cost share purchased from each domestic sector); the imported-input cost
share vector $\bm\Omega^F$; the labor cost share $\bm\alpha$; household consumption shares
$\bm\beta^H$; and each sector's export intensity (exports as a share of own gross output), used in
the sector-specific export-reallocation robustness check of Section \ref{sec:robustness}.

Table \ref{tab:chile_calib} reports the resulting Chile calibration --- the parameterization used for
every result in Sections \ref{sec:results} and \ref{sec:robustness} unless stated otherwise --- along
with the derived network objects (Domar weights $\lambda_{D,i}$, total import centrality $M_i$, and
the divine-coincidence weights $w_i^{DC}$) used throughout the paper.

\begin{table}[H]
\centering
\caption{Chile calibration: primitives and derived network objects}
\label{tab:chile_calib}
\begin{tabular}{lccc}
\toprule
 & Resource (1) & Manufacturing (2) & Services (3) \\
\midrule
Labor share $\alpha_i$ & 0.5026 & 0.3589 & 0.6035 \\
Import share $\Omega^F_i$ & 0.0767 & 0.1945 & 0.0704 \\
Consumption share $\beta^H_i$ & 0.0265 & 0.2294 & 0.7441 \\
Export share of own output & 0.6019 & 0.1797 & 0.0357 \\
Calvo reset probability $\delta_i$ & 0.90 & 0.31 & 0.16 \\
\midrule
Domar weight $\lambda_{D,i}$ & 0.0718 & 0.3820 & 1.1084 \\
Total import centrality $M_i$ & 0.1547 & 0.2847 & 0.1189 \\
Effective Calvo persistence $\hat\delta_i$ & 0.8902 & 0.1246 & 0.0311 \\
DC index weight $w_i^{DC}$ & 0.0002 & 0.0720 & 0.9277 \\
\bottomrule
\end{tabular}
\end{table}

\noindent Domestic IO matrix (rows buy from columns, Resource/Manufacturing/Services order):
$\Omega^H = \begin{psmallmatrix}0.0750 & 0.1526 & 0.1932\\ 0.0991 & 0.2022 & 0.1453 \\ 0.0018 & 0.0581 &
0.2661\end{psmallmatrix}$.

\subsection{Calvo Frequencies}

Sector-level Calvo reset probabilities $\delta_i$ are not identifiable from input--output data alone
and are taken from the price-stickiness literature (Dhyne et al., 2005, the ECB's Inflation
Persistence Network; cross-checked against Nakamura and Steinsson, 2008), assigning Resource the
highest reset probability (0.90, reflecting globally-priced commodity-like goods), Manufacturing an
intermediate value (0.31), and Services the lowest (0.16, reflecting the well-documented empirical
stickiness of service prices). The resulting effective Calvo persistence $\hat\delta_i$ (Table
\ref{tab:chile_calib}) is markedly lower for Services (0.031) than for Resource (0.890): a one
percent shock to Services' marginal cost translates into roughly thirty times less inflation, and
correspondingly more absorbed into markup dispersion, than the same shock in the Resource sector ---
the origin of Services' outsized welfare weight $\kappa_i$ discussed in Section \ref{sec:results}.

\subsection{Preference and Technology Parameters}

The remaining parameters --- the coefficient of relative risk aversion $\gamma=1$, the inverse Frisch
elasticity of labor supply $\varphi=2$, the within-sector elasticity of substitution
$\varepsilon=8$, the discount factor $\beta=0.99$, the home-foreign consumption elasticity
$\eta=1.5$, the export-demand elasticity $\theta^*=2.0$, and the debt-elasticity coefficient of the
risk premium $\Psi=0.020$ --- are literature defaults rather than estimates specific to this setting,
and are flagged as such in the discussion of the paper's remaining limitations (Section
\ref{sec:discussion}). All exogenous shocks (sectoral productivity, import price, foreign demand,
export price/terms-of-trade, and the risk premium itself) are calibrated to a common one percent
standard deviation, in the absence of an external target for the risk-premium shock specifically ---
a calibration choice whose consequences for the headline result are the subject of extensive
robustness analysis in Section \ref{sec:robustness}.

%======================================================================
\section{Quantitative Results}
\label{sec:results}
%======================================================================

\subsection{Headline Welfare Ranking}

Table \ref{tab:headline} reports the model's central quantitative result: welfare loss, in units of
$10^{-4}$ (i.e., basis points of steady-state consumption-equivalent loss times 100), for each of the
three monetary/exchange-rate regimes, computed from equation \eqref{eq:welfare_main}'s output-gap-plus-
within-sector terms using the order-1 (log-linear) solution's unconditional variances.

\begin{table}[H]
\centering
\caption{Headline welfare loss by regime, real Chile calibration ($\times10^{-4}$)}
\label{tab:headline}
\begin{tabular}{lccc}
\toprule
 & Float & Managed Float & Peg \\
\midrule
Welfare loss ($\times10^{-4}$) & 25.47 & 10.17 & 102.05 \\
Relative to Managed & 2.5$\times$ & 1$\times$ (best) & 10.0$\times$ \\
\bottomrule
\end{tabular}
\end{table}

Managed float dominates both corner regimes by a wide margin; Peg is roughly four times worse than
Float and ten times worse than Managed float. This is, superficially, the same ranking the classical
Mundell--Fleming/Gal\'i--Monacelli literature would predict --- flexible exchange rates are
stabilizing, hard pegs are costly --- but Section \ref{sec:results}'s shock decomposition shows the
mechanism generating this ranking is not the one that literature emphasizes.

Including the cross-sector welfare term of equation \eqref{eq:welfare_main} (Section \ref{sec:theory},
Remark on Cobb--Douglas cross-sector aggregators) raises every regime's loss modestly ---
3.3\% for Float, 8.4\% for Managed, and 4.5\% for Peg --- to totals of 26.31, 11.03, and 106.62
($\times10^{-4}$) respectively, without affecting the ranking; unless stated otherwise, all
subsequent tables in this paper use the output-gap-plus-within-sector total for direct comparability
with the extensive robustness campaign of Section \ref{sec:robustness}, none of which recomputes the
cross-sector term at every grid point.

\subsection{Shock Decomposition: The Risk-Premium Channel}
\label{sec:shock_decomp}

Table \ref{tab:shock_decomp} decomposes each regime's total welfare loss by the shock responsible for
it, using Dynare's own forecast-error variance decomposition. The result is stark: the risk-premium
shock $\varepsilon^{RP}_t$ --- a purely financial, non-import shock entering only through the UIP
condition \eqref{eq:uip_main} --- accounts for 75\% of Peg's total welfare loss at the model's
baseline calibration, far exceeding the terms-of-trade (export-price) channel that is the star of the
classical open-economy literature (a mere 2.16 out of 102.05, or roughly 2\%, of Peg's total loss).

\begin{table}[H]
\centering
\caption{Share of each regime's welfare loss attributable to the risk-premium shock ($\%$)}
\label{tab:shock_decomp}
\begin{tabular}{lccc}
\toprule
 & Float & Managed & Peg \\
\midrule
Risk-premium shock's \% of total loss & $\approx$11 & $\approx$8 & \textbf{75} \\
\bottomrule
\end{tabular}
\end{table}

This finding is quantitatively central to the paper, and uncomfortable in one specific sense that I
confront directly: Section \ref{sec:robustness}'s risk-premium volatility sweep shows that with the
risk-premium shock switched off entirely, Peg is not dominated at all --- it is essentially tied with
Float, fractionally \emph{better}. The entire ``Peg is worst'' headline result therefore depends on
the risk-premium shock existing at roughly its calibrated size, a calibration choice (equal standard
deviation to every other shock, for lack of an independently estimated target) rather than an
externally disciplined one. I return to this in Section \ref{sec:robustness}, where I show the
qualitative ranking survives even substantial variation in both the shock's size and the closing
device's debt-elasticity coefficient, while being explicit that the \emph{margin} of Peg's dominance
is highly sensitive to both.

\subsection{Why Services? A Sectoral Mechanism Decomposition}
\label{sec:why_services}

A natural puzzle, raised directly in early feedback on this project, is why the Services sector ---
the least import-exposed sector in the calibration (Table \ref{tab:chile_calib}, lowest direct
$\Omega^F_i$) and, on that basis alone, the sector one might expect to be most insulated from
exchange-rate-driven welfare costs --- in fact absorbs the largest share of welfare loss in every
regime. Three candidate explanations were raised, and all three turn out to be real and compounding
rather than competing.

\textbf{Indirect exposure via the network.} Total import centrality $M_i$ (equation
\eqref{eq:gamma_def}'s numerator, up to normalization) splits into a direct part ($\Omega^F_i$, imports
the sector buys itself) and an indirect part inherited via domestic suppliers. For Services, 41\% of
its total import centrality is indirect --- almost as large a share as Resource's 50\%, despite
Services' total exposure being the smallest of the three sectors in levels. Concretely, Services buys
5.8\% of its cost from Manufacturing, the most import-intensive sector in the economy; whatever
exchange-rate cost-push hits Manufacturing directly is inherited by Services through this single
supply link, and because Services cannot reprice quickly (its effective Calvo persistence
$\hat\delta_3=0.031$ is an order of magnitude lower than Resource's), whatever cost-push it inherits
this way does not clear within a quarter --- it lingers, making Services' inflation the single most
persistent series in the model (autocorrelation around 0.82).

\textbf{Which channel actually dominates Services' own variance?} Dynare's own variance decomposition
of Services' output gap shows the import-price/network channel just described is real but
\emph{secondary} --- explaining only 2--22\% of Services' output-gap variance depending on the regime
--- while the purely aggregate, non-import risk-premium shock dominates: 63\% under Float, 78\% under
Peg.

\textbf{Why does a non-import shock land so heavily on Services specifically?} Because of its Domar
weight, $\lambda_{D,3}=1.108$ --- more than double the next-largest sector's --- which reconciles the
apparent tension: 92\% of this weight is already present at zero network density, simply from
Services' large household consumption share and heavy within-sector input reuse; the network itself
adds only an additional 9\%. Resource, by contrast, sees its (much smaller) Domar weight more than
double when the network is switched on. The general lesson, applicable well beyond this specific
calibration: size (the Domar weight $\lambda_{D,i}$) determines which sector absorbs a generic
aggregate shock; rigidity ($\hat\delta_i$) converts that exposure into welfare cost; and the network
amplifies an already-large sector's exposure more than it creates new exposure from nothing.

\subsection{Which Sector Prefers Which Regime}

A related and initially counterintuitive finding, developed formally in Appendix
\ref{app:proofs} \S\ref{app:sector_pref}, is that the sector most directly exposed to the exchange
rate on impact (Resource, with the largest exchange-rate pass-through coefficient $\Gamma_1=0.234$
versus Services' near-zero $\Gamma_3=0.006$) is \emph{not} the sector that most prefers a peg. Pegging
does not eliminate the need for relative-price adjustment following a foreign-price or productivity
shock --- it relocates that adjustment into domestic (home-goods) prices, which, by construction, are
dominated by the upstream sector. This relocated-adjustment cost for Resource under a peg exceeds the
direct exchange-rate cost-push it avoids, so Resource, despite having the largest $\Gamma_i$, nets out
\emph{preferring float} between the two corner regimes. Services, three links removed from the
border and the stickiest sector, is doubly insulated from the peg's relocated-adjustment burden;
its inflation variance falls by nearly half under a peg relative to float, and because Services'
welfare weight $\kappa_3$ dwarfs the other two sectors', its own preference for the peg (conditional
on the two corner regimes only) dominates the aggregate ranking between Float and Peg. Managed float
dominates for \emph{every} sector simultaneously, precisely because it damps the exchange rate's
contribution to inflation variance without fully re-imposing the peg's relative-price-relocation
cost --- it need not pick a sectoral winner to improve on both corner regimes at once.

\subsection{Cross-Country Robustness: Korea and Czechia}
\label{sec:crosscountry}

Table \ref{tab:crosscountry} repeats the headline welfare comparison using independently constructed
national IO calibrations for South Korea and Czechia. The ranking --- Managed float dominant, Peg
worst by a wide margin --- replicates in both, with both economies' losses uniformly higher than
Chile's, consistent with both having denser domestic networks and higher import exposure than Chile
in the underlying IO data --- itself supporting evidence that the mechanism is a genuine network
story rather than a Chile-specific artifact.

\begin{table}[H]
\centering
\caption{Cross-country welfare loss by regime ($\times10^{-4}$)}
\label{tab:crosscountry}
\begin{tabular}{lccc}
\toprule
 & Float & Managed Float & Peg \\
\midrule
Chile & 25.47 & 10.17 & 102.05 \\
South Korea & 38.91 & 14.43 & 133.81 \\
Czechia & 32.20 & 11.82 & 84.03 \\
\bottomrule
\end{tabular}
\end{table}

A separate, weaker test asks whether Services' ``mostly indirect exposure via the network'' finding
(Section \ref{sec:why_services}) is a general pattern across downstream, low-import sectors. Pooling
all three calibrations (nine sector-level observations) gives only a moderate, noisy relationship
between how much a sector sources domestically and how much of its import exposure is indirect --- an
honest limitation given only three sectors per calibration. What holds robustly across all three
calibrations \emph{separately}, however, is that Services sources the \emph{least} domestically of any
sector in every single calibration, yet still derives a large share of its exposure indirectly,
because the small amount it does source domestically is concentrated in the single most import-heavy
supplier rather than spread evenly --- a composition effect, not a volume effect. A fully
disaggregated, many-sector calibration (Section \ref{sec:discussion}) is the natural next step to test
this pattern with statistical power beyond nine observations.

%======================================================================
\section{Robustness Analysis}
\label{sec:robustness}
%======================================================================

Because the headline welfare ranking rests, per Section \ref{sec:shock_decomp}, so heavily on two
specific calibration choices --- the risk-premium shock's standard deviation $\sigma_{RP}$ and the
debt-elasticity coefficient $\Psi$ --- this section subjects the model to an extensive robustness
campaign: in total, roughly 350 additional Dynare solves beyond the three headline country
calibrations, organized into a sequence of exercises. Throughout, all real-Chile-calibration exercises
hold the sector count, Calvo frequencies, and all shock processes other than the one(s) explicitly
varied fixed at their Table \ref{tab:chile_calib} baseline values, and network-density exercises use a
continuous scaling dial $\rho\in[0,2]$ applied to the off-diagonal entries of $\Omega^H$ ($\rho=0$: no
domestic cross-sector network, each sector uses only labor and imports; $\rho=1$: the real Chile
calibration; $\rho=2$: double density), with the diagonal (own-sector reuse) held fixed and the labor
share $\alpha_i$ absorbing the freed or additionally consumed cost share so that cost shares continue
to sum to one at every $\rho$.

\subsection{Network Density and the Optimal Managed-Float Response}

Does the network itself matter for welfare, beyond simply being ``an open economy''? Scaling network
density $\rho$ continuously on a stylized triangular network calibrated to match the real Chile
economy's total off-diagonal cost-share mass, the network premium --- the gap between welfare loss
with the network at its real density versus none at all --- roughly triples the Float/Peg welfare gap
as density rises, and reallocates \emph{which} sector bears the cost by regime. Does the
\emph{optimal} managed-float FX-response weight $\phi_s^*$ itself shift with network density, or is
the network-density result purely a level effect? A twelve-point grid search over $\phi_s$ at each of
three network densities ($\rho=0,1,2$) shows the optimum shifts materially: $\phi_s^*=0.15$ at
$\rho=0$, rising to $0.20$ at the real Chile density, and to $0.30$ at double density --- confirming
that denser production networks call for a \emph{more} aggressive exchange-rate response, not just a
larger welfare loss at any given response. The welfare curve is flat within roughly 2--3\% of each
optimum across a wide neighborhood of $\phi_s$, so this is a real, monotonic shift rather than a
knife-edge result, but getting $\phi_s$ within about $\pm0.1$ of correct matters far more than getting
it exactly right.

\subsection{Risk-Premium Volatility}
\label{rob:rp_vol}

Scaling the risk-premium shock's standard deviation from zero to three times its calibrated value,
holding every other shock and parameter fixed, shows that with the shock switched off entirely, Peg is
not dominated at all --- it is fractionally \emph{better} than Float (a ratio of 0.95). Half of the
calibrated shock size already flips this to Peg being 1.7 times worse than Float; the full calibrated
size gives the headline 3.6$\times$; and three times the calibrated size gives 14.5$\times$. This is
the single most important robustness finding in the paper in the sense that it directly qualifies the
headline result: the ``Peg is dominated'' conclusion is not an artifact of an implausible shock size
(even a modest fraction of the calibrated size already produces the qualitative ranking), but its
\emph{margin} is highly sensitive to a shock size that is, honestly, an assumed-equal-to-every-other-
shock convention rather than an externally estimated moment.

\subsection{$\psi$-Sensitivity, on Its Own and Jointly with Network Density and Shock Size}
\label{rob:psi}

The debt-elastic risk-premium coefficient $\Psi$ (baseline 0.020) is itself a literature default with
a known numerical pitfall: values too close to zero can generate a near-unit-root, knife-edge blowup
in net-foreign-asset dynamics that is a linearization artifact rather than real economics. Scaling
$\Psi$ from 0.25$\times$ to 4$\times$ its baseline (0.005 to 0.080) shows the qualitative ranking (Peg
worst) is robust across the entire grid, but the \emph{margin} is not constant: the Peg/Float welfare
ratio shrinks from 5.4$\times$ at the lowest $\Psi$ to 2.7$\times$ at the highest, and the risk-premium
shock's share of Peg's total loss (the 75\% headline figure, at baseline $\Psi=0.020$) ranges from 83\%
at the lowest $\Psi$ down to 58\% at the highest --- risk-premium/UIP remains the single dominant
channel throughout, but the exact 75\% figure is itself a function of an under-scrutinized calibration
choice.

Crossing this $\psi$-sensitivity check with network density (45 additional solves: five $\Psi$ values
$\times$ three densities $\times$ three regimes) shows the pattern found on the real Chile calibration
generalizes essentially unchanged: the Peg/Float ratio moves from roughly 5.3--5.5$\times$ at the
lowest $\Psi$ to 2.7--2.8$\times$ at the highest, at \emph{every} network density from zero to double
--- $\psi$-sensitivity and network density are essentially orthogonal margins of the model.

A separate, nine-point joint grid (scaling $\Psi$ and the risk-premium shock's own volatility
$\sigma_{RP}$ simultaneously, Peg regime only) tests whether the two calibration choices flagged above
reinforce or offset one another. A log-log regression of Peg's welfare loss on both scale factors
finds an elasticity with respect to $\sigma_{RP}$ of approximately $+1.44$ (dominant and convex ---
doubling the shock's size more than doubles the loss) against an elasticity with respect to $\Psi$ of
approximately $-0.24$ (real but secondary), with a non-trivial interaction term of approximately
$-0.16$: $\Psi$'s dampening effect on Peg's loss is \emph{stronger}, not weaker, when the risk-premium
shock itself is larger (the ratio $W(\Psi{=}2\times)/W(\Psi{=}0.5\times)$ shrinks from 0.85 at low
$\sigma_{RP}$ to 0.62 at high $\sigma_{RP}$). The two calibration choices flagged in Section
\ref{sec:shock_decomp} reinforce, rather than offset, one another.

\subsection{Price Rigidity: Uniform and Services-Specific}
\label{rob:rigidity}

Two complementary cuts test whether the ``Why Services?'' mechanism of Section \ref{sec:why_services}
is specifically about \emph{Services'} own rigidity, or reflects aggregate rigidity more generally.

\textbf{Uniform rigidity.} Scaling all three sectors' stickiness $(1-\hat\delta_i)$ by a common
multiplier $\kappa\in\{0.5,0.75,1.0,1.15\}$, crossed with network density, is bounded above by a real
feasibility constraint: the current calibration's Services sector is already close enough to the
Calvo-rigidity ceiling that a multiplier of 1.5 or above implies a \emph{negative} implied reset
probability, caught cleanly by the sweep script's own feasibility guard rather than by a numerical
crash. Within the feasible range, Peg's welfare loss rises monotonically and steeply in aggregate
rigidity at every network density (for example, at double network density: 60, 104, 182, and 251,
in units of $10^{-4}$, across the feasible $\kappa$ grid) --- rigidity compounds with Peg's existing
risk-premium problem, a clean and unambiguous result. Float and Managed, by contrast, are
\emph{not} monotonic across the same grid, and the topmost feasible point ($\kappa=1.15$) should be
read as a stress test rather than a reliable point: at that point, Services' Calvo persistence index
$\hat\delta_3$ collapses to 0.0015, an order of magnitude closer to zero than the already-small
baseline value (0.031), right at the model's own numerical feasibility floor, and the welfare weight
in equation \eqref{eq:welfare_main} scales like $(1-\hat\delta_i)/\hat\delta_i$, making this point
numerically fragile by construction.

\textbf{Services-only rigidity.} Moving only Services' own Calvo reset probability
$\delta_3\in\{0.05,0.10,0.16(\text{baseline}),0.25,0.40,0.60\}$, holding both other sectors fixed at
their baseline calibration, produces a cleaner picture, with no near-degenerate points anywhere in the
grid: Peg's welfare loss falls \emph{monotonically} as Services becomes more flexible, at every
network density (for example at the real Chile density: 122.9, 113.7, 102.1, 85.0, 60.7, and 37.0, in
units of $10^{-4}$, as $\delta_3$ rises from 0.05 to 0.60) --- a clean, robust confirmation that
Services' own rigidity is a first-order driver of Peg's welfare dominance, not merely a coincidental
correlate of aggregate stickiness. Float and Managed, in contrast, are hump-shaped rather than
monotonic in this same grid, peaking around $\delta_3\approx0.25$--$0.40$: the welfare weight
$(1-\hat\delta_i)/\hat\delta_i$ and the Phillips-curve slope's own sensitivity to $\delta_3$ move in
offsetting directions as rigidity falls, a genuine (if secondary) non-monotonicity rather than a
numerical artifact, since this grid stays well clear of the feasibility floor identified in the
uniform-rigidity cut above.

\subsection{Network Topology: Beyond the Density Dial}
\label{rob:topology}

Every exercise up to this point varies only the \emph{density} of a single, fixed triangular network
shape. This subsection asks whether the results are specific to that shape by constructing two
alternative topologies at matched total off-diagonal cost-share mass (0.6501, identical to the
baseline triangular network): a \textbf{hub-and-spoke} network with Manufacturing as the central node
(Resource and Services each link only to Manufacturing, with no direct Resource--Services link), and a
\textbf{chain} network with strictly one-directional pass-through, Resource $\to$ Manufacturing $\to$
Services, and no back-links at all.

The welfare ranking --- Peg worst, Managed best --- survives in \emph{every} topology (Table
\ref{tab:topology}), and the chain topology, with its most concentrated pass-through structure, is
uniformly worse than the baseline triangle for every regime, with the hub-and-spoke topology
intermediate. The optimal managed-float FX-response weight shifts with topology as well as with
density: $\phi_s^*=0.20$ for the triangle, but $0.30$ for both the hub-and-spoke and chain topologies,
consistent with (and reinforcing) the density-based finding of Section \ref{rob:rp_vol}'s companion
exercise above. Recomputing the analytical import-centrality decomposition of Section
\ref{sec:why_services} for each topology shows Services' indirect-exposure share is not an artifact of
the triangular shape --- if anything it strengthens under more concentrated pass-through structures,
rising from 41\% (triangle) to 46\% (hub-and-spoke) to 63\% (chain) of Services' total import
centrality.

\begin{table}[H]
\centering
\caption{Welfare loss by network topology and regime ($\times10^{-4}$)}
\label{tab:topology}
\begin{tabular}{lccc}
\toprule
Topology & Float & Managed & Peg \\
\midrule
Triangle (baseline) & 25.5 & 10.2 & 102.1 \\
Hub-and-spoke (Manufacturing hub) & 28.9 & 11.3 & 113.5 \\
Chain (Resource$\to$Manuf.$\to$Services) & 33.7 & 12.7 & 115.5 \\
\bottomrule
\end{tabular}
\end{table}

\subsection{Regime-Rule Variants}
\label{rob:regime_variants}

The baseline float regime targets Rubbo's divine-coincidence index with conventional (non-aggressive)
Taylor coefficients $(\phi_\pi,\phi_y)=(1.5,0.5)$. This subsection asks two further questions: does
targeting a different price index (strict domestic/PPI or strict consumer-price/CPI inflation,
instead of the divine-coincidence index) matter, and does the Taylor rule's own aggressiveness matter?

At the baseline, non-aggressive Taylor coefficients, adding two new policy-rule variants --- strict
CPI-inflation targeting and strict PPI (domestic-bundle)-inflation targeting, in place of the
divine-coincidence index --- to the model (verified, via a regression check re-running the three
original regimes through the edited model files, to reproduce the original headline numbers exactly,
confirming the new branches are additive and non-breaking) shows both new variants land statistically
indistinguishable from Float: 25.36 (CPI) and 25.26 (PPI) versus Float's 25.47 ($\times10^{-4}$), all
within about one percent of one another. \textbf{Targeting the theoretically privileged
divine-coincidence index, at these conventional policy coefficients, buys almost nothing over naively
targeting CPI or PPI inflation}; essentially all of Managed float's $\approx$2.5$\times$ improvement
comes from its explicit exchange-rate response term, not from smarter inflation targeting.

This picture changes once the Taylor rule is made aggressive. Setting $\phi_\pi=5$ (a numerical proxy
for a very aggressive, near-strict inflation-targeting rule) and $\phi_y=0$, at the real Chile network
density, gives a markedly different ranking among the three price-index targets: strict PPI targeting
(10.55, $\times10^{-4}$) outperforms strict divine-coincidence targeting (12.57), which in turn
outperforms strict CPI targeting (18.48). \textbf{Once policy is aggressive enough, targeting the
divine-coincidence index is no longer the best choice among simple alternatives; targeting domestic
(PPI) inflation strictly is}, reversing the conventional-coefficient finding above. A companion coarse
grid search over $(\phi_\pi,\phi_y)\in\{1.5,3,5\}\times\{0,0.5,1\}$ on the divine-coincidence-index
(float) rule finds its best point in this $3\times3$ grid, $(\phi_\pi,\phi_y)=(5,0.5)$, gives a welfare
loss of 9.55 ($\times10^{-4}$) --- essentially matching, and fractionally beating, Managed float's
headline 10.17, \emph{without any explicit exchange-rate response term at all}. I flag this finding
explicitly as preliminary rather than settled: it is a coarse $3\times3$ grid on the float branch only,
and Managed float's own $(\phi_\pi,\phi_y)$ coefficients were left at their original, non-aggressive
baseline rather than re-optimized at the same level of aggressiveness, so the comparison is not yet
apples-to-apples. What it does establish cleanly is that the baseline $\phi_\pi=1.5$ materially
undersells how well plain divine-coincidence-index targeting can do; $\phi_\pi$ itself had never
previously been swept in this project.

\subsection{The Network Premium on the Real Calibration}
\label{rob:network_premium}

Every network-density exercise above uses a stylized triangular network to permit continuous scaling.
As a direct check that the network channel matters on the paper's \emph{actual} headline calibration
--- not merely on an illustrative proxy --- I construct a no-network counterfactual by zeroing all
nine entries of Chile's real, dense $\Omega^H$ matrix (including the diagonal, so each sector uses
only labor and imports), letting the labor share $\alpha_i$ absorb the freed cost share, re-deriving
the non-stochastic steady state, and re-simulating all three regimes. The resulting network premium is
substantial and consistent across regimes: welfare loss with the real network in place, relative to
the no-network counterfactual, is 213\% higher for Float (25.47 versus 8.13), 169\% higher for Peg
(102.05 versus 37.94), and 92\% higher for Managed (10.17 versus 5.30). The production network itself
--- not simply the fact of being an open economy --- drives the majority of the welfare loss in every
regime on the paper's real, headline calibration, and amplifies Float's loss proportionally more than
Managed's.

\subsection{Sector-Specific Export Demand}
\label{rob:export}

The baseline model allocates aggregate exports across sectors via household consumption shares
$\bm\beta^H$ (heavily weighted toward Services), rather than via each sector's actual, empirically
very different export intensity (heavily weighted toward Resource: 60\% of Resource's own output is
exported, versus under 4\% for Services). This understates the model's ability to represent
``the export-heavy sector is hit by a terms-of-trade shock; how does the rest of the economy feel it
through the network,'' since under the baseline allocation the shock always spreads via consumption
shares rather than genuine export exposure. Replacing the single aggregate export equation with three
sector-specific export equations, recalibrated (via a damped fixed-point iteration on each sector's
export-scale constant) so that steady-state export intensity matches the real, empirical export shares
while keeping the aggregate trade balance and net-foreign-asset target unchanged, and re-deriving the
steady state accordingly, shows the export-price (terms-of-trade) shock's contribution to Peg's
welfare loss \emph{falls} substantially once exports are allocated realistically: from 2.16 to 1.17
($\times10^{-4}$), a 46\% reduction, because the baseline consumption-share allocation was overstating
how much export-price risk the (Services-heavy) consumption-share proxy exposes Peg to, relative to
the true (Resource-heavy) export exposure. Peg's total welfare loss falls correspondingly, from 102.05
to 90.45 ($\times10^{-4}$), an 11.4\% reduction, though the qualitative ranking is entirely unaffected.
A companion network-premium check under this more realistic export allocation reproduces the Section
\ref{rob:network_premium} finding closely (Float +190\%, Peg +169\%, Managed +88\%), and additional
sensitivity sweeps over the export-demand elasticity and the reallocation intensity both show smooth,
monotonic behavior with no ranking flips.

\subsection{Second-Order (Nonlinear) Welfare Check}
\label{rob:second_order}

All results above use the model's order-1 (log-linear) solution, evaluating the (already
second-order) welfare functional \eqref{eq:welfare_main} at first-order-accurate equilibrium paths ---
the standard linear-quadratic shortcut, valid whenever the first-order solution is certainty-equivalent.
This equivalence can, in principle, fail in this specific model because of two open-economy features
absent from Rubbo's closed economy: the debt-elastic UIP wedge is nonlinear in the net foreign asset
position, and the net foreign asset position itself has a root close to the inverse discount factor,
so precautionary (Jensen's-inequality) terms need not average out over any short horizon. Re-solving
the model to a genuine order-2 (pruned) approximation and simulating 260,000 periods (rather than
relying on Dynare's analytic higher-order moments, which require stationarity the model's own
price-level and exchange-rate-level state variables do not satisfy under an inflation-targeting rule),
computing $\mathbb E[X^2]$ directly from the simulated path rather than from the sample variance, and
comparing against an identical-seed, identical-length order-1 simulation to control for Monte Carlo
noise, shows the ranking is fully second-order-robust: every regime's loss rises by at most 2\% (Float
+2.0\%, Managed +0.6\%, Peg +0.7\%), nowhere near enough to threaten Managed $<$ Float $<$ Peg.
Decomposing the correction further shows the risk-adjusted mean shifts implied by the nonlinear model
are real and economically interpretable --- Peg's risk-adjusted output gap sits roughly ten times
larger, in absolute value, than Float's or Managed's, and net foreign assets carry a larger
precautionary buffer exactly where the UIP/risk-premium channel is least offset by exchange-rate
movement --- but because the welfare loss is quadratic, squaring these small mean shifts contributes
less than half a percent of the total squared loss in every regime; essentially all of the modest
second-order correction instead reflects a small, common variance amplification across all three
regimes, not a risk-adjusted-steady-state effect specific to any one of them.

\subsection{Summary of the Robustness Campaign}

Table \ref{tab:robustness_summary} collects the qualitative conclusion of every exercise in this
section. The headline ranking Managed $<$ Float $\ll$ Peg is, without a single exception across
roughly 350 additional model solves spanning network density, network topology, uniform and
sector-specific price rigidity, the risk-premium shock's size and closing-device elasticity (on their
own and jointly, and jointly with network density), and a range of alternative policy-rule designs,
qualitatively robust. What is \emph{not} robust to the same degree is the precise \emph{margin} of
Peg's dominance, which is highly sensitive to the risk-premium shock's calibrated volatility and its
interaction with the debt-elasticity coefficient $\Psi$ --- exactly the two parameters this paper's
underlying financial-account machinery introduces relative to Rubbo's closed economy and to Qiu, Wang,
Xu and Zanetti's (2026) static open economy.

\begin{table}[H]
\centering
\small
\caption{Robustness campaign summary: does Managed $<$ Float $\ll$ Peg survive?}
\label{tab:robustness_summary}
\begin{tabular}{p{4.3cm}lp{6.3cm}}
\toprule
Exercise & Ranking survives? & Key additional finding \\
\midrule
Network density ($\rho\in[0,2]$) & Yes & Network premium roughly triples Float/Peg gap \\
Optimal $\phi_s$ vs.\ density & --- & $\phi_s^*$: 0.15 $\to$ 0.20 $\to$ 0.30 as $\rho$ rises \\
Risk-premium volatility & Yes (mostly) & Off entirely: Peg $\approx$ tied with Float \\
$\psi$ sensitivity & Yes & Peg/Float ratio: 5.4$\times \to$ 2.7$\times$ across grid \\
$\psi\times$ network density & Yes & Interaction is negligible \\
$\psi\times\sigma_{RP}$ joint & Yes & Genuine reinforcing interaction ($\approx-0.16$) \\
Uniform rigidity $\times$ density & Yes & Peg monotonic; Float/Managed not \\
Services-only rigidity $\times$ density & Yes & Peg monotonic; Float/Managed hump-shaped \\
Network topology & Yes & Chain worst; $\phi_s^*$ shifts with shape too \\
Strict CPI/PPI-IT (baseline coeff.) & Yes & All three $\approx$ indistinguishable from Float \\
Strict CPI/PPI-IT (aggressive coeff.) & Yes & PPI $<$ DC-index $<$ CPI (reversed!) \\
Dual-mandate $(\phi_\pi,\phi_y)$ grid & Yes (caveated) & Aggressive plain float may rival Managed \\
Network premium (real calibration) & --- & +213\%/+169\%/+92\% (Float/Peg/Managed) \\
Sector-specific exports & Yes & Peg's loss falls 11.4\% under realistic exports \\
Second-order (nonlinear) welfare & Yes & Losses rise $\le$2\% in every regime \\
Cross-country (Korea, Czechia) & Yes & Both losses uniformly higher than Chile's \\
\bottomrule
\end{tabular}
\end{table}

%======================================================================
\section{Discussion}
\label{sec:discussion}
%======================================================================

\subsection{Policy Implications}

Three practical implications follow from this paper's results, with appropriate caveats attached to
each. First, for a small open economy with a production network resembling this paper's calibrations,
a managed float with a modest, positive response to the exchange rate is likely to dominate both
corner regimes by a wide margin, and this conclusion appears robust to essentially every dimension of
model misspecification examined in Section \ref{sec:robustness}. Second, the specific mechanism
generating this ranking --- a financial risk-premium/UIP channel, not the terms-of-trade channel
emphasized by the classical literature --- means that policies aimed at reducing exposure to the
risk-premium channel specifically (for example, reserve accumulation along the lines of Fanelli and
Straub, 2021, or policies reducing the debt-elasticity of the country's risk premium) may be a more
direct lever on welfare than exchange-rate policy design per se, a possibility this paper's model is
not currently equipped to evaluate directly since $\Psi$ and $\sigma_{RP}$ are treated as exogenous
primitives rather than policy-influenced objects. Third, Section \ref{rob:regime_variants}'s finding
that policy aggressiveness, not merely the choice of price index, is a first-order determinant of
welfare --- and that the theoretically privileged divine-coincidence index is not obviously the best
practical target once policy is aggressive --- suggests that calibrating how hard a central bank
\emph{can} lean against inflation, given its other constraints, may matter as much as getting the
target index exactly right.

\subsection{Limitations and Directions for Future Work}

Honestly assessed, this paper is not yet ready for submission to a top-five journal, and treating the
current draft as a checkpoint rather than a near-final product is the right frame. Several concrete
gaps remain. First and most importantly, the three-sector aggregation used throughout, while
sufficient to establish the qualitative mechanism and to permit the extensive robustness campaign of
Section \ref{sec:robustness}, is coarser than a competing paper in this space (Qiu, Wang, Xu and
Zanetti, 2026) that has already cleared peer review using a considerably more disaggregated
cross-border input--output dataset (the World Input-Output Database, 43 countries by 56 sectors); a
genuinely multi-sector calibration using an analogous disaggregated database (for example, OECD's
Trade in Value Added tables) is the most important remaining step, and is likely to be requested by
referees. Second, and directly related to the robustness campaign's central finding, several
parameters driving the headline result --- most importantly the risk-premium shock's standard
deviation and the debt-elasticity coefficient $\Psi$, but also the within- and across-sector
elasticities of substitution and the shock persistence parameters --- are literature defaults rather
than estimates specific to this setting or this country; an independently estimated risk-premium
process, in particular, is a necessary input before the margin (though not, per Section
\ref{sec:robustness}, the qualitative ranking) of the peg-dominated result can be treated as more than
suggestive. Third, the paper currently contains no empirical validation section tying the model's
quantitative magnitudes to data from an actual historical exchange-rate regime episode; Silva's (2024)
Chilean CPI-elasticity application is a natural template for such an exercise, and a natural
complement rather than a substitute for this paper's own normative, regime-comparison exercise. Fourth,
even a complete, referee-ready draft addressing the above typically undergoes two to three referee
rounds over one to two or more years at the journals this paper aspires to; that clock has not yet
started. The necessary condition for a genuine contribution --- the paper fills a gap explicitly
flagged as open by the authors of the closest competing paper themselves (Section
\ref{sec:litreview}) --- appears to be met; the sufficient condition, a complete and fully
referee-proofed paper, remains realistically several months to a year of further work away, beginning
with the multi-sector calibration and a fuller engagement with Qiu, Wang, Xu and Zanetti's (2026)
output-gap-weight derivation directly in the model section rather than solely as a literature-review
citation.

Beyond these gaps required for submission readiness, several extensions raised during the robustness
campaign (Section \ref{sec:robustness}) but flagged there as preliminary would repay further work: a
formally re-optimized managed-float rule at the same level of policy aggressiveness as the dual-mandate
grid search of Section \ref{rob:regime_variants}, to settle whether an aggressive plain float rule
genuinely rivals (rather than merely appears to rival) the managed float once both are given the same
optimization effort; and a dominant-currency-pricing extension along the lines of Gopinath and
Itskhoki (2022), replacing this paper's producer-currency-pricing assumption with a mixed invoicing
regime, which would alter the exchange-rate pass-through vector $\Gamma$ in ways this paper's current
model cannot represent.

%======================================================================
\section{Conclusion}
\label{sec:conclusion}
%======================================================================

This paper extends Rubbo's (2024) closed-economy theory of production networks and monetary policy to
a small open economy, and shows that her central divine-coincidence result --- the existence of a
single inflation index whose targeting simultaneously closes the output gap --- generically breaks
down once the exchange rate is introduced as a second, independent cost-push channel alongside
sectoral productivity. Quantitatively, calibrating the model to three real national input--output
datasets, I find that a managed float dominates both a free float and a hard peg, and that this
ranking is driven overwhelmingly by a financial risk-premium/uncovered-interest-parity channel rather
than the terms-of-trade mechanism emphasized by the classical open-economy literature. An extensive
robustness campaign --- spanning network density and shape, price rigidity, the calibration of the
risk-premium process, and alternative monetary policy rules --- confirms this qualitative ranking is
remarkably stable, while surfacing two further, more nuanced findings about policy aggressiveness and
price-index choice that merit further, more careful investigation before being treated as settled.
The paper's relationship to the contemporaneous and independent work of Qiu, Wang, Xu and Zanetti
(2026) is, I have argued, one of corroboration and natural extension rather than redundancy: their own
stated top-priority direction for future research --- relaxing the assumption of financial autarky ---
is precisely the machinery this paper's contribution is built around, and precisely the channel this
paper's quantitative analysis identifies as the true driver of the exchange-rate-regime welfare
ranking.

%======================================================================
\newpage
\section*{References}
\addcontentsline{toc}{section}{References}
%======================================================================

\begin{list}{}{\leftmargin=1.5em \itemindent=-1.5em \itemsep=0.5em}

\item Acemoglu, D., Carvalho, V. M., Ozdaglar, A. and Tahbaz-Salehi, A. (2012). The Network Origins of
Aggregate Fluctuations. \textit{Econometrica}, 80(5), 1977--2016.

\item Amiti, M., Itskhoki, O. and Konings, J. (2014). Importers, Exporters, and Exchange Rate
Disconnect. \textit{American Economic Review}, 104(7), 1942--1978.

\item Auer, R., Levchenko, A. A. and Sauré, P. (2019). International Inflation Spillovers Through
Input-Output Linkages. \textit{Review of Economics and Statistics}, 101(3), 507--521.

\item Baqaee, D. R. and Farhi, E. (2019). The Macroeconomic Impact of Microeconomic Shocks: Beyond
Hulten's Theorem. \textit{Econometrica}, 87(4), 1155--1203.

\item Baqaee, D. R. and Farhi, E. (2020). Productivity and Misallocation in General Equilibrium.
\textit{Quarterly Journal of Economics}, 135(1), 105--163.

\item Bems, R. and Johnson, R. C. (2017). Demand for Value Added and Value-Added Exchange Rates.
\textit{American Economic Journal: Macroeconomics}, 9(4), 45--90.

\item Bigio, S. and La'O, J. (2020). Distortions in Production Networks. \textit{Review of Economic
Studies}, 87(5), 2212--2246.

\item Calvo, G. A. (1983). Staggered Prices in a Utility-Maximizing Framework. \textit{Journal of
Monetary Economics}, 12(3), 383--398.

\item Carvalho, C. (2006). Heterogeneity in Price Stickiness and the Real Effects of Monetary Shocks.
\textit{B.E. Journal of Macroeconomics (Frontiers)}, 6(3).

\item Carvalho, V. M. and Tahbaz-Salehi, A. (2019). Production Networks: A Primer. \textit{Annual
Review of Economics}, 11, 635--663.

\item Corsetti, G., Dedola, L. and Leduc, S. (2010). Optimal Monetary Policy in Open Economies. In
B. M. Friedman and M. Woodford (eds.), \textit{Handbook of Monetary Economics}, Vol. 3B. Elsevier.

\item Dhyne, E., Álvarez, L. J., Le Bihan, H., Veronese, G., Dias, D., Hoffmann, J., Jonker, N., Lünnemann,
P., Rumler, F. and Vilmunen, J. (2005). Price Setting in the Euro Area: Some Stylized Facts from
Individual Consumer Price Data. European Central Bank Working Paper No.\ 524.

\item di Giovanni, J., Levchenko, A. A. and Méjean, I. (2018). The Micro Origins of International
Business Cycle Comovement. \textit{American Economic Review}, 108(1), 82--108.

\item Fanelli, S. and Straub, L. (2021). A Theory of Foreign Exchange Interventions. \textit{Review of
Economic Studies}, 88(6), 2857--2885.

\item Galí, J. and Monacelli, T. (2005). Monetary Policy and Exchange Rate Volatility in a Small Open
Economy. \textit{Review of Economic Studies}, 72(3), 707--734.

\item Ghassibe, M. (2021). Monetary Policy and Production Networks: An Empirical Investigation.
\textit{Journal of Monetary Economics}, 119, 21--39.

\item Gopinath, G. and Itskhoki, O. (2022). Dominant Currency Paradigm: A Review. In G. Gopinath, E.
Helpman and K. Rogoff (eds.), \textit{Handbook of International Economics}, Vol. 6. Elsevier.

\item Hulten, C. R. (1978). Growth Accounting with Intermediate Inputs. \textit{Review of Economic
Studies}, 45(3), 511--518.

\item La'O, J. and Tahbaz-Salehi, A. (2022). Optimal Monetary Policy in Production Networks.
\textit{Econometrica}, 90(3), 1295--1336.

\item Nakamura, E. and Steinsson, J. (2008). Five Facts about Prices: A Reevaluation of Menu Cost
Models. \textit{Quarterly Journal of Economics}, 123(4), 1415--1464.

\item Obstfeld, M. and Rogoff, K. (2000). New Directions for Stochastic Open Economy Models.
\textit{Journal of International Economics}, 50(1), 117--153.

\item Ozdagli, A. and Weber, M. (2017). Monetary Policy Through Production Networks. NBER Working
Paper No.\ 23764.

\item Pasten, E., Schoenle, R. and Weber, M. (2020). The Propagation of Monetary Policy Shocks in a
Heterogeneous Production Economy. \textit{Journal of Monetary Economics}, 116, 37--52.

\item Qiu, Z., Wang, Y., Xu, L. and Zanetti, F. (2026). Monetary Policy in Open Economies with
Production Networks. \textit{Journal of Monetary Economics}, 159, 103918.

\item Rubbo, E. (2024). Networks, Phillips Curves, and Monetary Policy. \textit{Econometrica}, 92(5),
1373--1408.

\item Schmitt-Grohé, S. and Uribe, M. (2003). Closing Small Open Economy Models. \textit{Journal of
International Economics}, 61(1), 163--185.

\item Schmitt-Grohé, S. and Uribe, M. (2016). Downward Nominal Wage Rigidity, Currency Pegs, and
Involuntary Unemployment. \textit{Journal of Political Economy}, 124(5), 1466--1514.

\item Silva, A. (2024). Inflation in Disaggregated Small Open Economies. Federal Reserve Bank of
Boston Working Paper 24-12; arXiv:2410.00705.

\end{list}

%======================================================================
\newpage
\appendix
\section{Proofs}
\label{app:proofs}
%======================================================================

This appendix collects detailed, step-by-step proofs of every lemma and proposition stated in Section
\ref{sec:theory}. The closed-economy results (Lemmas \ref{lem:natural_output}--\ref{app:lem2} and
Propositions \ref{prop:pc_closed}--\ref{prop:dc_closed}) reproduce Rubbo's (2024) own derivations in
full, restated here for completeness and because the open-economy results
(\S\ref{app:open_economy_proofs}) are obtained by an essentially identical argument applied to a
modified marginal cost equation.

\subsection{Notation}

Lower-case letters denote log-deviations from a zero-inflation, zero-markup non-stochastic steady
state throughout. $\bm\alpha\in\mathbb R^N$ is the labor cost-share vector, $\Omega^H\in\mathbb
R^{N\times N}$ the domestic input--output matrix, $\bm\beta^H\in\mathbb R^N$ household consumption
shares across domestic sectors, $\Delta=\mathrm{diag}(\hat\delta_1,\ldots,\hat\delta_N)$ the effective
Calvo persistence matrix, and $(I-\Omega^H)^{-1}$, $(I-\Omega^H\Delta)^{-1}$ the (rigidity-unadjusted
and rigidity-adjusted) Leontief inverses. $\bm\lambda^\top=(\bm\beta^H)^\top(I-\Omega^H)^{-1}$ are
Domar weights.

\subsection{Log-Linearized Building Blocks}

\textbf{Marginal cost.} Since production is constant-returns-to-scale and productivity is
Hicks-neutral, the unit cost function is $MC_{it}=\tilde c_i(W_t,\{P_{jt}\})/A_{it}$. By Shephard's
lemma, $\partial\log MC_{it}/\partial\log W_t=\alpha_i$ (labor's cost share) and
$\partial\log MC_{it}/\partial\log P_{jt}=\Omega^H_{ij}$ (input $j$'s cost share); and
$\partial\log MC_{it}/\partial\log A_{it}=-1$. Log-linearizing and stacking across sectors,
\begin{equation}
    \bm{mc}_t = \bm\alpha\,w_t + \Omega^H\bm p_t - \log\mathbf A_t. \label{app:eq:mc}
\end{equation}

\textbf{Price-level dynamics and Calvo reset price.} With a fraction $\delta_i$ of sector-$i$ firms
resetting price each period, log-linearizing $P_{it}^{1-\varepsilon_i}=(1-\delta_i)P_{it-1}^{1-\varepsilon_i}
+\delta_i(P_{it}^*)^{1-\varepsilon_i}$ around the zero-inflation steady state gives
$p_{it}^*-p_{it-1}=\pi_{it}/\delta_i$. Log-linearizing the optimal Calvo reset price around the same
steady state gives $p_{it}^*=(1-\rho(1-\delta_i))\sum_{s\ge0}[\rho(1-\delta_i)]^s\mathbb E_tmc_{it+s}$.
Combining these two and solving forward one step yields the sectoral Phillips curve
\begin{equation}
    \pi_{it} = \hat\delta_i(mc_{it}-p_{it-1}) + \rho(1-\hat\delta_i)\mathbb E_t\pi_{it+1},
    \qquad
    \hat\delta_i \equiv \frac{\delta_i(1-\rho(1-\delta_i))}{1-\rho\delta_i(1-\delta_i)},
    \label{app:eq:pilaw}
\end{equation}
or, in matrix form, $\bm\pi_t=\Delta(\bm{mc}_t-\bm p_{t-1})+\rho(I-\Delta)\mathbb E\bm\pi_{t+1}$.

\textbf{Real wage.} Log-linearizing labor supply $W_t/P_t^C=C_t^\gamma L_t^\varphi$ and using goods
market clearing $c_t\approx y_t$ together with Lemma \ref{lem:natural_output} below,
\begin{equation}
    w_t - (\bm\beta^H)^\top\bm p_t = (\gamma+\varphi)\tilde y_t + \bm\lambda^\top\log\mathbf A_t.
    \label{app:eq:realwage}
\end{equation}

\subsection{Proof of Lemma \ref{lem:natural_output} (Natural Output)}

The flexible-price equilibrium is efficient, so it solves the planner's problem
$\max\, U(\mathcal C(\{Y_i-\sum_jX_{ji}\}),L)$ subject to $Y_i=A_iF_i(\{X_{ij}\},L_i)$. Decompose this
into an inner problem --- given total labor $L$, allocate it (and produce) efficiently across sectors,
defining $C^*(L;\mathbf A)$ --- and an outer problem choosing $L$. By constant returns to scale of
every $F_i$ and of the consumption aggregator $\mathcal C$, $C^*$ is homogeneous of degree one in $L$.
The outer first-order condition is $C^{*-\gamma}\partial C^*/\partial L = L^\varphi$. By the envelope
theorem applied to the inner Lagrangian, $\partial C^*/\partial L=\nu_L$ (the shadow value of labor)
and $\partial C^*/\partial A_i = \nu_i\,\partial(A_iF_i)/\partial A_i = \nu_i Y_i^*$ (the shadow value
of good $i$ times its optimal output). Converting to logs and using that, in competitive equilibrium,
$\nu_i=P_i/(P^CC^{*\gamma})$, gives $\partial\log C^*/\partial\log A_i = P_iY_i^*/(P^CC^*)\equiv
\lambda_i$, the Domar weight. Because $C^*$ is homogeneous of degree one in $L$, Euler's theorem gives
$C^*=L\nu_L$; combining this with the outer first-order condition and the labor-supply and
labor-demand conditions in the natural equilibrium (using $y_{\mathrm{nat}}-l_{\mathrm{nat}}=
\bm\lambda^\top\log\mathbf A_t$, itself following from Step 4's Domar-weight identity applied to the
aggregate CRS production function) yields, after collecting terms,
$y_{t,\mathrm{nat}}=\tfrac{1+\varphi}{\gamma+\varphi}\bm\lambda^\top\log\mathbf A_t$. $\blacksquare$

\subsection{Proof of Lemma (Output Gap Equals Employment Gap)}
\label{app:lem2}

By constant returns to scale, aggregate output in both the actual and the natural equilibrium is
given by the same inner value function, $y_t=C_t^*(L_t;\mathbf A_t)$,
$y_{t,\mathrm{nat}}=C_t^*(L_{t,\mathrm{nat}};\mathbf A_t)$. Log-linearizing around $L_{t,\mathrm{nat}}$
and applying the envelope theorem exactly as in the proof of Lemma \ref{lem:natural_output},
$\partial\log C^*/\partial\log L = L\nu_L/C^*$, which equals one at the natural equilibrium because
$\nu_L=C^{*\gamma}L^\varphi$ and the outer first-order condition holds with equality. Hence
$\tilde y_t\equiv y_t-y_{t,\mathrm{nat}} = 1\cdot(l_t-l_{t,\mathrm{nat}})\equiv\tilde l_t$.
$\blacksquare$

\subsection{Proof of Lemma (Output Gap Equals Minus the Sales-Weighted Markup)}

From the household's labor-supply condition, in both the actual and natural equilibria,
$(w_t-p_t^C)-(w_{t,\mathrm{nat}}-p_{t,\mathrm{nat}}^C)=(\gamma+\varphi)\tilde y_t$, using
$\tilde y_t=\tilde l_t$ from the preceding lemma. Writing the price gap
$\tilde{\bm p}_t\equiv\bm p_t-\bm p_{t,\mathrm{nat}}$, the marginal-cost gap satisfies
$\tilde{\bm{mc}}_t=\bm\alpha\tilde w_t+\Omega^H\tilde{\bm p}_t$ (subtracting equation
\eqref{app:eq:mc} evaluated at the natural equilibrium, where $\log\mathbf A_t$ cancels). Since markups
are $\mu_{it}=p_{it}-mc_{it}$ and are zero in the natural equilibrium,
$\tilde{\bm p}_t=\tilde{\bm{mc}}_t+\bm\mu_t=\bm\alpha\tilde w_t+\Omega^H\tilde{\bm p}_t+\bm\mu_t$, so
$(I-\Omega^H)\tilde{\bm p}_t=\bm\alpha\tilde w_t+\bm\mu_t$, and pre-multiplying by
$(I-\Omega^H)^{-1}$ and then by $(\bm\beta^H)^\top$ gives $(\bm\beta^H)^\top\tilde{\bm p}_t =
(\bm\lambda^\top\bm\alpha)\tilde w_t + \bm\lambda^\top\bm\mu_t = \tilde w_t + \bm\lambda^\top\bm\mu_t$,
using the constant-returns-to-scale identity $\bm\lambda^\top\bm\alpha=1$ (equivalently
$(I-\Omega^H)^{-1}\bm\alpha=\mathbf 1$, since $\alpha_i+\sum_j\Omega^H_{ij}=1$ for every $i$).
Substituting into the labor-supply-gap equation above,
$\tilde w_t - (\tilde w_t+\bm\lambda^\top\bm\mu_t)=(\gamma+\varphi)\tilde y_t$, i.e.
$(\gamma+\varphi)\tilde y_t=-\bm\lambda^\top\bm\mu_t$. $\blacksquare$

\subsection{Proof of Proposition \ref{prop:pc_closed} (Sectoral Phillips Curve)}

Substituting equation \eqref{app:eq:mc} into \eqref{app:eq:pilaw} and using $\bm p_t=\bm
p_{t-1}+\bm\pi_t$ gives, after collecting all $\bm\pi_t$ terms on the left,
$(I-\Delta\Omega^H-\Delta\bm\alpha(\bm\beta^H)^\top)\bm\pi_t = \Delta\bm\alpha(\gamma+\varphi)\tilde
y_t + \Delta(\bm\alpha\bm\lambda^\top-I)\log\mathbf A_t + \Delta(\bm\alpha(\bm\beta^H)^\top-I+\Omega^H)
\bm p_{t-1} + \rho(I-\Delta)\mathbb E\bm\pi_{t+1}$, where the real-wage equation
\eqref{app:eq:realwage} has been substituted for $w_t$. The left-hand matrix factors as
$(I-\Delta\Omega^H)(I-\Delta(I-\Omega^H\Delta)^{-1}\bm\alpha(\bm\beta^H)^\top)$, using the identity
$(I-\Delta\Omega^H)^{-1}\Delta=\Delta(I-\Omega^H\Delta)^{-1}$ (verified by direct multiplication). The
second factor is a rank-one perturbation of the identity, invertible by the Sherman--Morrison formula,
$(I-\mathbf u\mathbf v^\top)^{-1}=I+\mathbf u\mathbf v^\top/(1-\mathbf v^\top\mathbf u)$, with
$\mathbf u=\Delta(I-\Omega^H\Delta)^{-1}\bm\alpha\equiv\tilde A\bm\alpha$ and $\mathbf v=\bm\beta^H$.
Defining $\kappa\equiv1-(\bm\beta^H)^\top\tilde A\bm\alpha$, direct computation of
$(I-\Delta\Omega^H-\Delta\bm\alpha(\bm\beta^H)^\top)^{-1}\Delta\bm\alpha$ using this inverse gives
exactly $\tilde A\bm\alpha/\kappa\equiv\mathcal B$, the output-gap coefficient in equation
\eqref{eq:pc_closed}. Carrying the analogous computation through for the productivity and lagged-price
terms, and collecting them into $-\mathcal V\bm\chi_t$ with
$\bm\chi_t\equiv(I-\Omega^H)^{-1}\log\mathbf A_t-\bm p_{t-1}$, yields equation \eqref{eq:pc_closed}.
That $\mathcal V\mathbf 1=\mathbf 0$ follows because $\mathcal V=(A-I)(I-\Omega^H)$ for a matrix $A$
satisfying $A\mathbf 1=\mathbf 1$ (verified directly from $A$'s definition using
$(I-\Omega^H)^{-1}\bm\alpha=\mathbf 1$), so $\mathcal V\mathbf 1=(A-I)(I-\Omega^H)\mathbf
1=(A-I)\bm\alpha=A\bm\alpha-\bm\alpha$, which vanishes because $A\bm\alpha=\bm\alpha$ by the same
identity applied componentwise. $\blacksquare$

\subsection{Proof of Proposition \ref{prop:dc_closed} (Divine Coincidence, Closed Economy)}

From equation \eqref{app:eq:pilaw} and the markup definition $\mu_{it}=p_{it}-mc_{it}$, substituting
$mc_{it}-p_{it-1}=(\pi_{it}-\rho(1-\hat\delta_i)\mathbb E_t\pi_{it+1})/\hat\delta_i$ gives the exact
markup-inflation link $\mu_{it}=-\tfrac{1-\hat\delta_i}{\hat\delta_i}(\pi_{it}-\rho\mathbb
E_t\pi_{it+1})$. Substituting this into $(\gamma+\varphi)\tilde y_t=-\bm\lambda^\top\bm\mu_t$ (the
preceding lemma) gives $(\gamma+\varphi)\tilde y_t = \big(\sum_i\lambda_i\tfrac{1-\hat\delta_i}
{\hat\delta_i}\big)(\pi_t^{DC}-\rho\mathbb E\pi_{t+1}^{DC})$ where $\pi_t^{DC}$ is exactly the
Domar-rigidity-weighted average defined in the proposition; rearranging gives the cost-push-free
Phillips curve \eqref{eq:DCPC} directly, with no productivity term surviving because only Lemma
\ref{lem:natural_output}--\ref{app:lem2}'s output-gap/markup link and the Calvo markup-inflation
identity were used, neither of which involves $\log\mathbf A_t$ directly.

For uniqueness: an index $\bm\phi^\top\bm\pi_t$ has no cost-push term if and only if
$\bm\phi^\top\mathcal V=\mathbf 0^\top$, i.e.\ $\bm\phi$ lies in $\mathcal V$'s left null space. Since
$\mathcal V\mathbf 1=\mathbf 0$ (established above), $\mathcal V$ has rank at most $N-1$, so its left
null space is generically exactly one-dimensional, and direct verification shows
$\bm\phi\propto\bm\lambda^\top(I-\Delta)\Delta^{-1}$ (component-wise, $\phi_i\propto
\lambda_i(1-\hat\delta_i)/\hat\delta_i$) satisfies $\bm\phi^\top\mathcal V=\mathbf 0^\top$. The
divine-coincidence weights are therefore the unique (up to normalization) solution. $\blacksquare$

\subsection{Proof of Proposition \ref{prop:pc_open} (Sectoral Phillips Curve, Open Economy)}
\label{app:open_economy_proofs}

Replace equation \eqref{app:eq:mc} with the open-economy marginal cost \eqref{eq:mc_main},
$\bm{mc}_t=\bm\alpha w_t+\Omega^H\bm p_t+\bm\Omega^F(p_t^*+e_t)-\log\mathbf A_t$. Every step of the
proof of Proposition \ref{prop:pc_closed} above goes through verbatim with this modified marginal
cost, since none of those steps used any property of $-\log\mathbf A_t$ beyond its being an additive,
sector-specific term entering marginal cost linearly --- exactly the property the new term
$\bm\Omega^F(p_t^*+e_t)$ shares. The new term is therefore pre-multiplied, in the final collected
equation, by exactly the same matrix, $(I-\Delta\Omega^H-\Delta\bm\alpha(\bm\beta^H)^\top)^{-1}\Delta$,
that pre-multiplies the output-gap coefficient, which was shown above to equal $\tilde A/\kappa$.
Since $p_t^*$ is exogenous (foreign, small-open-economy assumption) and only $e_t$ is a domestic
state variable, this delivers exactly the new term $\Gamma\Delta e_t$ in equation \eqref{eq:pc_open},
with $\Gamma=\tilde A\bm\Omega^F/\kappa$ as claimed, and every other term in equation
\eqref{eq:pc_open} unchanged from the closed economy. $\blacksquare$

\subsection{Proof of Proposition \ref{prop:dc_breakdown} (Generic Breakdown of the Divine Coincidence)}

Pre-multiplying equation \eqref{eq:pc_open} by a candidate weighting vector $\bm\phi^\top$ gives
$\pi_t^{\bm\phi} = \rho\bm\phi^\top(I-\mathcal V)\mathbb E\bm\pi_{t+1} + \bm\phi^\top\mathcal
B(\gamma+\varphi)\tilde y_t - \bm\phi^\top\mathcal V\bm\chi_t + (\bm\phi^\top\Gamma)\Delta e_t$; this
has no cost-push term of either kind if and only if both $\bm\phi^\top\mathcal V=\mathbf 0^\top$
\emph{and} $\bm\phi^\top\Gamma=0$ hold. By Proposition \ref{prop:dc_closed}, the first condition alone
already pins $\bm\phi$ down uniquely, up to scale, as the closed-economy divine-coincidence weight
vector $\bm\phi^{DC}$ --- there is no remaining freedom in $\bm\phi$ once the first condition is
imposed. The second condition, $(\bm\phi^{DC})^\top\Gamma=0$, is then a single additional scalar
equation with zero remaining free parameters to satisfy it: generically (for almost every choice of
the primitives $\Omega^H,\bm\alpha,\bm\beta^H,\Delta,\bm\Omega^F$, in the sense of Lebesgue measure on
the space of feasible primitives), this equation fails. It holds only in the non-generic case where
$\Gamma$ is itself proportional to $\mathcal B$ (equivalently, since $\mathcal B\propto\tilde
A\bm\alpha$ and $\Gamma\propto\tilde A\bm\Omega^F$, the case $\bm\Omega^F\propto\bm\alpha$: import
exposure exactly proportional to labor intensity in every sector), since in that knife-edge case
$\bm\phi^{DC\top}\Gamma\propto\bm\phi^{DC\top}\mathcal B$, and it can be verified directly from the
definitions that $\bm\phi^{DC\top}\mathcal B$ need not itself vanish, but the \emph{output-gap}
coefficient in the DC Phillips curve absorbs the entire exchange-rate term identically to how it
already absorbs the output-gap term, restoring a (modified) divine coincidence with a shifted
natural-rate object rather than eliminating the exchange-rate channel outright. $\blacksquare$

\subsection{Proof of Proposition \ref{prop:welfare} (Welfare)}

To second order, the representative household's utility loss relative to the efficient allocation
decomposes additively into: (i) an output-gap term, from the second-order expansion of
$U(C_t,L_t)$ around the efficient point using $\tilde y_t=\tilde l_t$ (the closed-economy Lemma
above), giving $-\tfrac{\gamma+\varphi}{2}U_C^*C^*\tilde y_t^2$; (ii) a within-sector term, from the
fact that under Calvo pricing the cross-sectional variance of log prices within sector $i$ is
$\mathrm{Var}_f[\log P_{if,t}]\approx\tfrac{1-\hat\delta_i}{\hat\delta_i}\pi_{it}^2$, which generates a
first-order-in-variance TFP loss of $-\tfrac12\varepsilon_i\mathrm{Var}_f[\log P_{if,t}]$ in sector
$i$, summed with Domar weights $\lambda_{D,i}$; and (iii) cross-sector terms, from the second-order
Baqaee--Farhi (2019) substitution operators $\Phi_C,\Phi_s$ evaluated at the model's consumption and
production elasticities of substitution, capturing the welfare cost of relative-price distortions
across sectors (as distinct from within-sector price dispersion). In the open economy, the marginal
cost equation \eqref{eq:mc_main} introduces an additional argument, $e_t\bm\Omega^F$, into these same
cross-sector operators, since different sectors face different import cost shocks and hence different
relative-price distortions even absent any within-sector dispersion; no new operators are needed, only
a new argument to the existing ones. $\blacksquare$

\begin{remark}[Closed form under Cobb--Douglas cross-sector nests]
\label{rem:cross_sector_numeric}
This model's cross-sector aggregators are Cobb--Douglas everywhere except the outer home/foreign
consumption nest: household consumption across domestic sectors, and each firm's intermediate-input
bundle across domestic sectors and imports, are both Cobb--Douglas by the assumed functional form of
$Y_{ift}$, and Cobb--Douglas delivers an Allen--Uzawa elasticity of substitution of exactly one between
any pair of its arguments, independent of the exponent (cost-share) weights. With a constant
elasticity $\sigma$ across all pairs, the algebraic identity $\sum_{k,h}w_kw_h(X_k-X_h)(Y_k-Y_h) =
2\big[\sum_kw_kX_kY_k - (\sum_kw_kX_k)(\sum_hw_hY_h)\big]$ (for weights summing to one) collapses
$\Phi_C(\bm\mu,\bm\mu)$ to a weighted variance of the markup-gap vector, with the import composite's
markup identically zero under the law of one price. Only the outer CES consumption nest ($\eta=1.5$)
contributes a genuinely non-unitary elasticity, entering as a practical stand-in for the true nested
cross-elasticity between domestic and imported goods. Numerically, in the Chile calibration, this
cross-sector term adds 3.3\% (Float), 8.4\% (Managed), and 4.5\% (Peg) on top of the output-gap-plus-
within-sector total, without affecting the regime ranking.
\end{remark}

\subsection{Which Sectors Prefer Which Regime: Supporting Derivation}
\label{app:sector_pref}

Two calibrated objects jointly determine each sector's regime preference. The first, direct
exchange-rate exposure, is $\Gamma_i=[\tilde A\bm\Omega^F]_i/\kappa$ from equation
\eqref{eq:gamma_def}: large when sector $i$ imports directly or buys heavily from suppliers who do,
discounted by how quickly prices reset along the chain. The second, the welfare weight from equation
\eqref{eq:welfare_main}, is $\kappa_i\equiv\lambda_{D,i}\varepsilon_i(1-\hat\delta_i)/\hat\delta_i$:
large when sector $i$ is a large, central domestic supplier and very sticky. In a triangular network
where the upstream sector is also the direct importer (as calibrated here, $\bm\Omega^F$ decreasing
downstream, $\bm\lambda_D$ increasing downstream), these two objects move in \emph{opposite}
directions across sectors: $\Gamma_1=0.234\gg\Gamma_3=0.006$, but $\kappa_1=0.43\ll\kappa_3=74.6$. The
sector most exposed to the exchange rate on impact is the one welfare cares about least per unit of
its inflation variance, and vice versa; neither parameter alone predicts the aggregate regime
ranking, but their product, combined with how each regime redistributes \emph{where} the necessary
relative-price adjustment occurs, does. Applying $\kappa_i$ to the simulated regime-by-regime
inflation variances gives each sector's realized loss $\kappa_i\mathrm{Var}^R(\pi_i)$: Sector 1's loss
is \emph{higher} under peg (0.59) than float (0.44), despite the $\Gamma$-channel switching off
entirely under a peg, because pegging relocates the adjustment burden into home-goods prices, which
sector 1 dominates; Sector 3's loss falls by nearly half under peg (16.96 under float versus 8.86
under peg) because it is both insulated from direct exchange-rate exposure and largely spared the
peg's relocated-adjustment burden. Since $\kappa_3\gg\kappa_1,\kappa_2$, Sector 3's preference for peg
over float dominates the aggregate ranking between the two corner regimes, even though Sector 1
individually prefers float. Managed float dominates for every sector simultaneously because a positive
$\phi_s$ damps the exchange rate's contribution to inflation variance without fully re-imposing the
peg's relocation cost, requiring no sectoral trade-off to still improve on both corner regimes.

\end{document}
